\chapter{Cohomology of sheaves of modules}
\label{chap:Cohomology}

Fix a commutative ring \(R\); for the genus of a curve, \(R\) will be the base field \(k\).
This chapter defines the cohomology of sheaves of \(R\)-modules on a small site as
\(\Ext\)-groups from the constant sheaf, and more generally the cohomology of an object of
the site as \(\Ext\)-groups from a free sheaf of \(R\)-modules; it identifies degree-zero
cohomology with (global) sections and proves the vanishing of degree-one cohomology of the
structure sheaf on an affine open of a scheme. It then develops the Mayer--Vietoris
\((0,1)\)-slice for a sheaf of \(R\)-modules on a site and specializes it to a scheme
covered by two affine opens, producing an explicit description of \(H^1(X,\struct{X/k})\)
as the cokernel of a map of section modules. It then runs the
two-lattice ladder of Chapter~\ref{chap:Algebra} on that cokernel along a finite morphism
to the projective line, proving the finiteness of \(H^1(C, \struct C)\) --- so that the
genus of Definition~\ref{def:genus} is the dimension of a finite-dimensional vector space.
The two closing sections glue a twisted sheaf out of a unit cocycle on a two-chart cover
and prove that the two-cover \(H^1\) of the relative curve --- for the structure sheaf,
and termwise for the twisted sheaf --- commutes with arbitrary change of the affine test
ring.

\section{Cohomology on a site}

Let \((\mathcal C, J)\) be a small site: a small category \(\mathcal C\) with a Grothendieck
topology \(J\). The category of sheaves of \(R\)-modules on \((\mathcal C, J)\) is a
Grothendieck abelian category; in particular it has enough injectives, and all
\(\Ext\)-groups exist and carry \(R\)-module structures.

\begin{definition}[The constant sheaf]
  \label{def:constModuleSheaf}
  \lean{CategoryTheory.Sheaf.constModuleSheaf}
  \leanok
  The constant sheaf \(R_J\) of \(R\)-modules on \((\mathcal C, J)\) with value \(R\): the
  sheafification of the constant presheaf with value \(R\).
\end{definition}

\begin{definition}[Sheaf cohomology]
  \label{def:HModule}
  \lean{CategoryTheory.Sheaf.HModule}
  \uses{def:constModuleSheaf}
  \leanok
  For a sheaf \(F\) of \(R\)-modules on \((\mathcal C, J)\) and \(n \ge 0\), the degree-\(n\)
  cohomology of \(F\) is the \(R\)-module
  \[
    H^n(F) \;=\; \Ext^n\bigl(R_J,\, F\bigr),
  \]
  the \(\Ext\)-group computed in the category of sheaves of \(R\)-modules.
\end{definition}

\begin{definition}[Functoriality]
  \label{def:HModule_map}
  \lean{CategoryTheory.Sheaf.HModule.map}
  \uses{def:HModule}
  \leanok
  A morphism of sheaves \(f : F \to G\) induces an \(R\)-linear map
  \(H^n(F) \to H^n(G)\) for every \(n\), by covariant functoriality of \(\Ext\) in the second
  argument. The identity induces the identity, and a composite induces the composite.
\end{definition}
\begin{proof}
  \leanok
  Functoriality of \(\Ext^n(R_J, -)\); the \(R\)-linearity comes from the \(R\)-linear
  structure of the sheaf category.
\end{proof}

\begin{lemma}[No higher cohomology on injectives]
  \label{lem:HModule_injective_vanishing}
  \lean{CategoryTheory.Sheaf.HModule.instSubsingletonHAddNatOfNat}
  \uses{def:HModule}
  \leanok
  If \(F\) is an injective object of the category of sheaves of \(R\)-modules, then
  \(H^{n+1}(F) = 0\) for every \(n \ge 0\).
\end{lemma}
\begin{proof}
  \leanok
  Positive-degree \(\Ext\)-groups into an injective object vanish.
\end{proof}

\begin{definition}[The adjunction, linearly]
  \label{def:constAdjHomEquiv}
  \lean{CategoryTheory.Sheaf.constantSheafAdjHomLinearEquiv}
  \leanok
  Assume the category \(\mathcal C\) has a terminal object \(\top\). For every \(R\)-module
  \(M\) and every sheaf \(F\) of \(R\)-modules there is an \(R\)-linear equivalence
  \[
    \Hom\bigl(\underline{M},\, F\bigr) \;\cong\; \Hom_R\bigl(M,\, F(\top)\bigr),
  \]
  where \(\underline{M}\) denotes the constant sheaf with value \(M\); it is the hom-bijection
  of the adjunction between the constant-sheaf functor and evaluation at \(\top\).
\end{definition}
\begin{proof}
  \leanok
  The constant-sheaf functor is left adjoint to evaluation at the terminal object. The
  bijection sends \(\varphi\) to the composite of the unit \(M \to \underline{M}(\top)\) with
  \(\varphi_\top\); both this assignment and its inverse are additive and commute with the
  \(R\)-action.
\end{proof}

\begin{definition}[Morphisms out of the constant sheaf are sections]
  \label{def:constHomEquiv}
  \lean{CategoryTheory.Sheaf.constModuleSheafHomEquiv}
  \uses{def:constModuleSheaf, def:constAdjHomEquiv}
  \leanok
  The \(R\)-linear equivalence
  \(\Hom(R_J, F) \cong F(\top)\): specialize \ref{def:constAdjHomEquiv} to \(M = R\) and
  compose with the canonical evaluation \(\Hom_R(R, N) \cong N\), \(u \mapsto u(1)\).
\end{definition}

\begin{lemma}[Naturality]
  \label{lem:constHomEquiv_naturality}
  \lean{CategoryTheory.Sheaf.constModuleSheafHomEquiv_naturality}
  \uses{def:constHomEquiv}
  \leanok
  The equivalence of \ref{def:constHomEquiv} is natural in the sheaf: for
  \(\varphi : R_J \to F\) and \(f : F \to G\), the section corresponding to
  \(f \circ \varphi\) is \(f_\top\) applied to the section corresponding to \(\varphi\).
\end{lemma}
\begin{proof}
  \leanok
  Naturality of the adjunction hom-bijection in the second variable, and naturality of the
  evaluation \(\Hom_R(R, N) \cong N\).
\end{proof}

\begin{definition}[Degree-zero cohomology is global sections]
  \label{def:HModule_linearEquiv0}
  \lean{CategoryTheory.Sheaf.HModule.linearEquiv₀}
  \uses{def:HModule, def:constHomEquiv}
  \leanok
  The \(R\)-linear equivalence \(H^0(F) \cong F(\top)\): the canonical identification
  \(\Ext^0(R_J, F) \cong \Hom(R_J, F)\) followed by \ref{def:constHomEquiv}.
\end{definition}

\begin{lemma}[Naturality in degree zero]
  \label{lem:HModule_linearEquiv0_naturality}
  \lean{CategoryTheory.Sheaf.HModule.linearEquiv₀_naturality}
  \uses{def:HModule_linearEquiv0, def:HModule_map}
  \leanok
  The equivalence \(H^0(F) \cong F(\top)\) of \ref{def:HModule_linearEquiv0} is natural: for
  \(f : F \to G\), applying \(f_\top\) to the section corresponding to \(x \in H^0(F)\)
  yields the section corresponding to the image of \(x\) under \(H^0(f)\).
\end{lemma}
\begin{proof}
  \leanok
  \uses{lem:constHomEquiv_naturality}
  The identification \(\Ext^0 \cong \Hom\) turns composition with the degree-zero class of
  \(f\) into postcomposition with \(f\); conclude by \ref{lem:constHomEquiv_naturality}.
\end{proof}

\section{Cohomology of an object of the site}

For an object \(U\) of \(\mathcal C\), this section defines the cohomology of \(U\) with
coefficients in a sheaf \(F\) of \(R\)-modules as \(\Ext\)-groups from a \emph{free} sheaf
of \(R\)-modules on \(U\), generalizing the previous section (where the implicit object was
the terminal object \(\top\)). This is the coefficient system needed to run the
Mayer--Vietoris sequence of Section~\ref{sec:mv-in-text} \(R\)-linearly, and it recovers the
cohomology of the site when \(U\) is terminal (Definition~\ref{def:HModule_linearEquivHModule'}
below).

\begin{definition}[The free sheaf of \(R\)-modules on an object of the site]
  \label{def:freeModuleSheaf}
  \lean{CategoryTheory.Sheaf.freeModuleSheaf}
  \leanok
  For an object \(U\) of \(\mathcal C\), the free sheaf of \(R\)-modules on \(U\), written
  \(R[U]\): the sheafification of the presheaf \(V \mapsto\) the free \(R\)-module on the
  set \(\Hom(V, U)\), with restriction maps induced by precomposition.
\end{definition}

\begin{definition}[Functoriality in the object of the site]
  \label{def:freeModuleSheafMap}
  \lean{CategoryTheory.Sheaf.freeModuleSheafMap}
  \uses{def:freeModuleSheaf}
  \leanok
  A morphism \(i : U \to V\) of \(\mathcal C\) induces a morphism of sheaves of
  \(R\)-modules \(R[i] : R[U] \to R[V]\): apply the free-\(R\)-module-on-a-set functor and
  sheafification to the induced map of representable presheaves
  \(\Hom(-,U) \to \Hom(-,V)\), \(g \mapsto i \circ g\).
\end{definition}

\begin{definition}[Free sheaves of modules represent sections]
  \label{def:freeModuleSheafHomEquiv}
  \lean{CategoryTheory.Sheaf.freeModuleSheafHomEquiv}
  \uses{def:freeModuleSheaf}
  \leanok
  For every object \(U\) of \(\mathcal C\) and every sheaf \(F\) of \(R\)-modules, an
  \(R\)-linear equivalence
  \[
    \Hom\bigl(R[U],\, F\bigr) \;\cong\; F(U).
  \]
\end{definition}
\begin{proof}
  \leanok
  Write \(y_U = \Hom(-,U)\) for the presheaf represented by \(U\), so that \(R[U]\) is the
  sheafification of the presheaf of \(R\)-modules obtained by applying the free-\(R\)-module
  functor to \(y_U\) objectwise. Compose three natural bijections:
  \begin{itemize}
    \item sheafification is left adjoint to the inclusion of sheaves into presheaves, giving
      \(\Hom(R[U],F) \cong \Hom_{\mathrm{PSh}}(y_U \otimes R,\, F)\), naturally in \(F\);
    \item on each object \(V\), the free-\(R\)-module functor is left adjoint to the
      forgetful functor from \(R\)-modules to sets, giving
      \(\Hom_{\mathrm{PSh}}(y_U \otimes R,\, F) \cong \Hom_{\mathrm{PSh}(\mathrm{Set})}(y_U,\, F)\)
      (\(F\) viewed as a presheaf of sets), naturally in \(F\);
    \item the Yoneda lemma identifies \(\Hom_{\mathrm{PSh}(\mathrm{Set})}(y_U, F) \cong F(U)\).
  \end{itemize}
  Each bijection is additive and \(R\)-linear in \(F\) (the first two are hom-sets of
  \(R\)-linear/additive adjunctions; the last carries the \(R\)-module structure of \(F(U)\)
  tautologically, evaluating a natural transformation at the identity of \(U\)), so the
  composite is an \(R\)-linear equivalence.
\end{proof}

\begin{lemma}[Naturality of the section equivalence]
  \label{lem:freeModuleSheafHomEquiv_naturality}
  \lean{CategoryTheory.Sheaf.freeModuleSheafHomEquiv_comp,
    CategoryTheory.Sheaf.freeModuleSheafHomEquiv_map}
  \uses{def:freeModuleSheafHomEquiv, def:freeModuleSheafMap}
  \leanok
  The equivalence of \ref{def:freeModuleSheafHomEquiv} is natural in both variables:
  \begin{enumerate}
    \item for \(\varphi : R[U] \to F\) and \(f : F \to G\), the section of \(G(U)\)
      corresponding to \(f \circ \varphi\) is \(f_U\) applied to the section of \(F(U)\)
      corresponding to \(\varphi\);
    \item for \(i : U \to V\) and \(\varphi : R[V] \to F\), the section of \(F(U)\)
      corresponding to the composite \(R[U] \xrightarrow{R[i]} R[V] \xrightarrow{\varphi} F\)
      is the restriction along \(i\) of the section of \(F(V)\) corresponding to \(\varphi\).
  \end{enumerate}
\end{lemma}
\begin{proof}
  \leanok
  Both are naturality of the composite of adjunction bijections used in
  \ref{def:freeModuleSheafHomEquiv}: (1) is naturality of each of the three bijections in
  the target sheaf \(F\); (2) is naturality of the Yoneda bijection
  \(\Hom_{\mathrm{PSh}(\mathrm{Set})}(y_U,-) \cong \mathrm{ev}_U\) in \(U\) (precomposition
  with \(y_i : y_U \to y_V\) corresponds to restriction along \(i\)), transported through the
  other two (fixed) bijections.
\end{proof}

\begin{definition}[Cohomology of an object of the site]
  \label{def:HModule'}
  \lean{CategoryTheory.Sheaf.HModule'}
  \uses{def:freeModuleSheaf}
  \leanok
  For an object \(U\) of \(\mathcal C\), a sheaf \(F\) of \(R\)-modules and \(n \ge 0\), the
  degree-\(n\) cohomology of \(U\) with coefficients in \(F\) is the \(R\)-module
  \[
    H'^n(U, F) \;=\; \Ext^n\bigl(R[U],\, F\bigr).
  \]
\end{definition}

\begin{definition}[Degree zero is sections]
  \label{def:HModule'_linearEquiv0}
  \lean{CategoryTheory.Sheaf.HModule'.linearEquiv₀}
  \uses{def:HModule', def:freeModuleSheafHomEquiv}
  \leanok
  The \(R\)-linear equivalence \(H'^0(U, F) \cong F(U)\): the canonical identification
  \(\Ext^0(R[U], F) \cong \Hom(R[U], F)\) followed by \ref{def:freeModuleSheafHomEquiv}.
\end{definition}

\begin{definition}[Restriction]
  \label{def:HModule'_res}
  \lean{CategoryTheory.Sheaf.HModule'.res}
  \uses{def:HModule', def:freeModuleSheafMap}
  \leanok
  For \(i : U \to V\) in \(\mathcal C\), restriction along \(i\) is the \(R\)-linear map
  \(H'^n(V, F) \to H'^n(U, F)\), \(x \mapsto x \circ R[i]\): precomposition with
  \(R[i] : R[U] \to R[V]\) (contravariant functoriality of \(\Ext^n(-, F)\)).
\end{definition}

\begin{definition}[Functoriality in the coefficient sheaf]
  \label{def:HModule'_mapCoeff}
  \lean{CategoryTheory.Sheaf.HModule'.mapCoeff,
    CategoryTheory.Sheaf.HModule'.mapCoeff_id_apply,
    CategoryTheory.Sheaf.HModule'.mapCoeff_comp_apply}
  \uses{def:HModule'}
  \leanok
  A morphism of sheaves of \(R\)-modules \(f : F \to G\) induces, for every object \(U\) of
  \(\mathcal C\) and every \(n\), an \(R\)-linear map \(H'^n(U, F) \to H'^n(U, G)\),
  \(x \mapsto x \cdot f\): composition (Yoneda product) with the degree-zero class of \(f\), i.e.\
  covariant functoriality of \(\Ext^n(R[U], -)\) in the coefficient variable. This is the
  coefficient-variable analogue of the restriction \ref{def:HModule'_res} (functoriality in the
  object of the site). The identity induces the identity, and a composite \(g \circ f\) induces the
  composite of the induced maps.
\end{definition}
\begin{proof}
  \leanok
  Postcomposition with the degree-zero \(\Ext\)-class of \(f\) is \(R\)-linear; it sends the
  degree-zero class of the identity to the identity operation, and by associativity of the Yoneda
  product it turns the class of a composite into the composite of the two postcompositions.
\end{proof}

\begin{definition}[Coefficient-isomorphism transport for cohomology of an object]
  \label{def:HModule'_congrCoeff}
  \lean{CategoryTheory.Sheaf.HModule'.congrCoeff}
  \uses{def:HModule'_mapCoeff}
  \leanok
  An isomorphism \(\alpha : F \cong G\) of coefficient sheaves of \(R\)-modules induces, for every
  object \(U\) and every \(n\), an \(R\)-linear isomorphism \(H'^n(U, F) \cong H'^n(U, G)\).
\end{definition}
\begin{proof}
  \leanok
  Take the induced maps of \ref{def:HModule'_mapCoeff} for \(\alpha\) and for \(\alpha^{-1}\). By
  functoriality their two composites are the maps induced by \(\alpha^{-1}\circ\alpha = \id_F\) and
  \(\alpha\circ\alpha^{-1} = \id_G\), which are the identity operations; hence the two induced maps
  are mutually inverse \(R\)-linear isomorphisms.
\end{proof}

\begin{lemma}[Transport of vanishing along a coefficient isomorphism]
  \label{lem:HModule'_subsingleton_of_congrCoeff}
  \lean{CategoryTheory.Sheaf.HModule'.subsingleton_of_congrCoeff}
  \uses{def:HModule'_congrCoeff}
  \leanok
  If \(F \cong G\) as coefficient sheaves of \(R\)-modules and \(H'^n(U, G) = 0\), then
  \(H'^n(U, F) = 0\).
\end{lemma}
\begin{proof}
  \leanok
  The \(R\)-linear isomorphism \(H'^n(U, F) \cong H'^n(U, G)\) of \ref{def:HModule'_congrCoeff}
  identifies \(H'^n(U, F)\) with the zero module \(H'^n(U, G)\).
\end{proof}

\begin{definition}[The free sheaf on a terminal object is the constant sheaf]
  \label{def:freeModuleSheafIsoConstModuleSheaf}
  \lean{CategoryTheory.Sheaf.freeModuleSheafIsoConstModuleSheaf}
  \uses{def:freeModuleSheaf, def:constModuleSheaf, def:freeModuleSheafHomEquiv,
    def:constHomEquiv}
  \leanok
  If \(T\) is a terminal object of \(\mathcal C\), a canonical isomorphism of sheaves of
  \(R\)-modules \(R[T] \cong R_J\).
\end{definition}
\begin{proof}
  \leanok
  Both sides corepresent the same functor: \(\Hom(R[T], -) \cong \mathrm{ev}_T\) by
  \ref{def:freeModuleSheafHomEquiv}, and \(\Hom(R_J, -) \cong \mathrm{ev}_T\) by
  \ref{def:constHomEquiv} (valid since \(T\) is terminal), both naturally in the sheaf
  argument. Two objects corepresenting the same functor via natural equivalences are
  canonically isomorphic: the isomorphism corresponds, under each equivalence, to the section
  of the other object at \(T\) that corresponds to its own identity morphism, and the two
  composites are the identity by naturality applied to the identity morphism of each side.
\end{proof}

\begin{definition}[Cohomology of the site is cohomology of a terminal object]
  \label{def:HModule_linearEquivHModule'}
  \lean{CategoryTheory.Sheaf.HModule.linearEquivHModule'}
  \uses{def:HModule, def:HModule', def:freeModuleSheafIsoConstModuleSheaf}
  \leanok
  For \(T\) terminal in \(\mathcal C\) and every \(n\), an \(R\)-linear equivalence
  \(H^n(F) \cong H'^n(T, F)\).
\end{definition}
\begin{proof}
  \leanok
  Apply the (contravariant, \(R\)-linear) functoriality of \(\Ext^n(-, F)\) to the
  isomorphism \(R[T] \cong R_J\) of \ref{def:freeModuleSheafIsoConstModuleSheaf}: an
  isomorphism of the first \(\Ext\)-argument induces, by composition (Yoneda product) with
  its degree-zero class, a natural \(R\)-linear isomorphism on \(\Ext\)-groups in every
  degree, with inverse given by composition with the inverse isomorphism.
\end{proof}

\section{The structure sheaf as a sheaf of modules}

Let \(k\) be a commutative ring and let \(X\) be a scheme over \(\Spec k\). The relevant site
is the \emph{small Zariski site} of \(X\): the partially ordered set of open subsets of
\(X\), with coverings the open coverings. The whole space is a terminal object of this site.

\begin{definition}[The algebra map on sections]
  \label{def:overAlgebraMap}
  \lean{AlgebraicGeometry.Scheme.overAlgebraMap}
  \leanok
  For every open \(U \subseteq X\), the canonical ring homomorphism \(k \to \Gamma(X, U)\):
  the map induced on global sections by the structure morphism, followed by restriction to
  \(U\). It commutes with the restriction maps of the structure sheaf; in particular each
  \(\Gamma(X, U)\) is a \(k\)-algebra --- hence a \(k\)-module --- and all restriction maps
  are \(k\)-linear.
\end{definition}
\begin{proof}
  \leanok
  Restriction maps are ring homomorphisms and commute with restriction from the global
  sections; the maps from \(k\) are fixed composites of such maps, so they commute with
  further restriction.
\end{proof}

\begin{definition}[The structure sheaf as a sheaf of \(k\)-modules]
  \label{def:moduleKSheaf}
  \lean{AlgebraicGeometry.Scheme.moduleKSheaf}
  \uses{def:overAlgebraMap}
  \leanok
  The structure sheaf of \(X\), viewed as a sheaf \(\struct{X/k}\) of \(k\)-modules on the
  small Zariski site of \(X\): its sections over an open \(U\) are \(\Gamma(X, U)\) with the
  \(k\)-module structure of \ref{def:overAlgebraMap}, and its restriction maps are those of
  the structure sheaf.
\end{definition}
\begin{proof}
  \leanok
  The restriction maps are \(k\)-linear by \ref{def:overAlgebraMap}, so this is a presheaf of
  \(k\)-modules. Its underlying presheaf of sets is that of the structure sheaf, which
  satisfies the sheaf condition; since limits of \(k\)-modules are computed on underlying
  sets, the presheaf of \(k\)-modules is a sheaf.
\end{proof}

\begin{definition}[Degree-zero cohomology of the structure sheaf]
  \label{def:moduleKSheafHZero}
  \lean{AlgebraicGeometry.Scheme.moduleKSheafHZero}
  \uses{def:moduleKSheaf, def:HModule_linearEquiv0}
  \leanok
  The \(k\)-linear equivalence \(H^0\bigl(X, \struct{X/k}\bigr) \cong \Gamma(X, \struct X)\).
\end{definition}
\begin{proof}
  \leanok
  The open \(X\) is a terminal object of the small Zariski site; apply
  \ref{def:HModule_linearEquiv0}.
\end{proof}

\section{\v{C}ech cobounding on affine opens}

For a scheme \(X\), a section \(g\) of the structure sheaf on an open \(V \subseteq X\) and
an open \(U \subseteq V\), we write \(D(g)\) for the basic open locus of \(V\) where \(g\) is
invertible, and \(g|_U\), \(s|_U\) for restrictions of sections.

\begin{lemma}[Denominator clearing on an affine open]
  \label{lem:denominator_clearing}
  \lean{AlgebraicGeometry.IsAffineOpen.exists_res_eq_pow_mul}
  \leanok
  Let \(X\) be a scheme, \(U \subseteq X\) an affine open, \(V \subseteq X\) an open with
  \(U \subseteq V\), and \(g \in \Gamma(X, V)\) a section of the structure sheaf on \(V\) (so
  that \(D(g) \subseteq V\) is its basic open). For every \(c \in \Gamma(X, U \cap D(g))\)
  there exist \(N \ge 0\) and \(e \in \Gamma(X, U)\) such that
  \[
    e|_{U \cap D(g)} \;=\; \bigl(g|_{U}\bigr)|_{U \cap D(g)}^{N} \cdot c .
  \]
\end{lemma}
\begin{proof}
  \leanok
  \(U \cap D(g)\) is the basic open of the affine \(U\) defined by \(g|_U\), and restriction
  identifies the sections over the basic open of an affine open with the localization:
  \(\Gamma(X, U \cap D(g)) = \Gamma(X, U)\bigl[1/(g|_U)\bigr]\) (quasi-coherence of the
  structure sheaf). Writing \(c\) as a fraction \(e / (g|_U)^N\) and clearing the denominator
  gives the claim.
\end{proof}

\begin{lemma}[Annihilation on an affine open]
  \label{lem:section_annihilation}
  \lean{AlgebraicGeometry.IsAffineOpen.exists_pow_mul_eq_zero_of_res_eq_zero}
  \leanok
  Let \(X\) be a scheme, \(U \subseteq X\) an affine open, \(V \subseteq X\) an open with
  \(U \subseteq V\), and \(g \in \Gamma(X, V)\) a section of the structure sheaf on \(V\) (so
  that \(D(g) \subseteq V\) is its basic open). If \(t \in \Gamma(X, U)\) restricts to zero
  on \(U \cap D(g)\), then \((g|_U)^{M} \cdot t = 0\) in \(\Gamma(X, U)\) for some \(M \ge 0\).
\end{lemma}
\begin{proof}
  \leanok
  As in \ref{lem:denominator_clearing}, \(\Gamma(X, U \cap D(g))\) is the localization of
  \(\Gamma(X, U)\) at the powers of \(g|_U\), and the kernel of the localization map consists
  of the elements annihilated by a power of \(g|_U\).
\end{proof}

\begin{theorem}[Degree-one \v{C}ech cobounding on an affine open]
  \label{thm:affine_open_cech_cobounding}
  \lean{AlgebraicGeometry.IsAffineOpen.exists_cech_cobounding}
  \leanok
  \source{stacks-project}
  Let \(X\) be a scheme and \(U\) an affine open of \(X\). Let \(f_1, \dots, f_n \in
  \Gamma(X, U)\) with \(D(f_1) \cup \dots \cup D(f_n) = U\). Let
  \(a_{ij} \in \Gamma\bigl(X, D(f_i) \cap D(f_j)\bigr)\) be a \v{C}ech \(1\)-cocycle: for all
  \(i, j, l\),
  \[
    a_{il} \;=\; a_{ij} + a_{jl}
    \qquad \text{on } D(f_i) \cap D(f_j) \cap D(f_l).
  \]
  Then there are sections \(b_i \in \Gamma\bigl(X, D(f_i)\bigr)\) with
  \[
    a_{ij} \;=\; b_i - b_j \qquad \text{on } D(f_i) \cap D(f_j)
  \]
  for all \(i, j\).
\end{theorem}
\begin{proof}
  \leanok
  \uses{lem:denominator_clearing, lem:section_annihilation}
  All restriction maps are suppressed from the notation. Note that each \(D(f_i)\) and each
  pairwise intersection \(D(f_i) \cap D(f_j) = D(f_i f_j)\) is an affine open of \(U\), hence
  of \(X\).

  \emph{Step 1: clear denominators with a uniform exponent.} For each pair \((i,j)\) apply
  \ref{lem:denominator_clearing} with affine open \(D(f_i)\), ambient open \(V = U\) and
  \(g = f_j\): there are \(N_{ij} \ge 0\) and \(e^0_{ij} \in \Gamma(X, D(f_i))\) with
  \(e^0_{ij} = f_j^{N_{ij}} a_{ij}\) on \(D(f_i) \cap D(f_j)\). Let \(N\) be the maximum of
  the \(N_{ij}\) and set \(e_{ij} = f_j^{\,N - N_{ij}}\, e^0_{ij}\), so that
  \(e_{ij} = f_j^{\,N} a_{ij}\) on \(D(f_i) \cap D(f_j)\).

  \emph{Step 2: kill the defects.} For \(i, j, l\) define, on \(D(f_i) \cap D(f_j)\),
  \[
    t_{ijl} \;=\; e_{il} - e_{jl} - f_l^{\,N} a_{ij}.
  \]
  On the triple intersection \(D(f_i) \cap D(f_j) \cap D(f_l)\) we have
  \(e_{il} = f_l^N a_{il}\) and \(e_{jl} = f_l^N a_{jl}\), so the cocycle identity gives
  \(t_{ijl} = f_l^N (a_{il} - a_{jl} - a_{ij}) = 0\) there. Since
  \(D(f_i) \cap D(f_j)\) is an affine open (of \(U\), hence of \(X\)) and the triple
  intersection is its basic open defined by \(f_l\), \ref{lem:section_annihilation} (with
  ambient open \(V = U\)) yields \(M_{ijl} \ge 0\) with
  \(f_l^{\,M_{ijl}}\, t_{ijl} = 0\) on \(D(f_i) \cap D(f_j)\). Let \(M\) be the maximum of
  the \(M_{ijl}\); then \(f_l^{\,M} t_{ijl} = 0\) for all \(i,j,l\).

  \emph{Step 3: partition of unity.} Since the \(D(f_i)\) cover the affine open \(U\), the
  \(f_i\) generate the unit ideal of \(\Gamma(X, U)\); hence so do the powers
  \(f_i^{\,N+M}\), and we may write \(1 = \sum_l h_l\, f_l^{\,N+M}\) with
  \(h_l \in \Gamma(X, U)\).

  \emph{Step 4: assemble the cobounding.} Set
  \(b_i = \sum_l h_l\, f_l^{\,M}\, e_{il} \in \Gamma\bigl(X, D(f_i)\bigr)\). On
  \(D(f_i) \cap D(f_j)\),
  \[
    b_i - b_j
      = \sum_l h_l f_l^{\,M} \bigl(e_{il} - e_{jl}\bigr)
      = \sum_l h_l f_l^{\,M} \bigl(f_l^{\,N} a_{ij} + t_{ijl}\bigr)
      = \Bigl(\sum_l h_l f_l^{\,N+M}\Bigr) a_{ij}
      = a_{ij},
  \]
  using \(f_l^{\,M} t_{ijl} = 0\) in the third equality.
\end{proof}

\begin{theorem}[Degree-one \v{C}ech cobounding on an affine scheme]
  \label{thm:cech_cobounding}
  \lean{AlgebraicGeometry.Scheme.exists_cech_cobounding}
  \uses{thm:affine_open_cech_cobounding}
  \leanok
  Let \(X\) be an affine scheme and let \(f_1, \dots, f_n\) be global sections of the
  structure sheaf with \(D(f_1) \cup \dots \cup D(f_n) = X\). Let
  \(a_{ij} \in \Gamma\bigl(X, D(f_i) \cap D(f_j)\bigr)\) be a \v{C}ech \(1\)-cocycle as in
  \ref{thm:affine_open_cech_cobounding}. Then there are sections
  \(b_i \in \Gamma\bigl(X, D(f_i)\bigr)\) with \(a_{ij} = b_i - b_j\) on
  \(D(f_i) \cap D(f_j)\) for all \(i, j\).
\end{theorem}
\begin{proof}
  \leanok
  The special case \(U = X\) of \ref{thm:affine_open_cech_cobounding}: an affine scheme is
  an affine open of itself.
\end{proof}

\section{Vanishing of degree-one cohomology on affine opens}

\begin{lemma}[A vanishing criterion for \(\Ext^1\)]
  \label{lem:ext_one_vanishing_criterion}
  \lean{CategoryTheory.Abelian.Ext.subsingleton_one_of_injective_of_surjective}
  \leanok
  Let \(\mathcal A\) be an abelian category (with \(\Ext\)-groups), let
  \(\iota : A \to I\) be a monomorphism with \(I\) injective, and let \(L\) be an object. If
  every morphism \(L \to \coker \iota\) lifts along the projection
  \(\pi : I \to \coker \iota\), then \(\Ext^1(L, A) = 0\).
\end{lemma}
\begin{proof}
  \leanok
  Consider the short exact sequence \(0 \to A \to I \to \coker \iota \to 0\) and the induced
  long exact sequence
  \[
    \Hom(L, I) \xrightarrow{\;\pi_*\;} \Hom(L, \coker \iota)
      \xrightarrow{\;\delta\;} \Ext^1(L, A) \longrightarrow \Ext^1(L, I).
  \]
  Let \(z \in \Ext^1(L, A)\). Its image in \(\Ext^1(L, I)\) vanishes because \(I\) is
  injective, so by exactness \(z = \delta(x)\) for some \(x : L \to \coker \iota\). By
  hypothesis \(x = \pi \circ \psi\) for some \(\psi : L \to I\), and
  \(\delta \circ \pi_* = 0\) (composite of consecutive maps of the long exact sequence);
  hence \(z = \delta(\pi_*(\psi)) = 0\).
\end{proof}

\begin{lemma}[Sections are left exact]
  \label{lem:sections_left_exact}
  \lean{CategoryTheory.Sheaf.exists_app_eq_of_cokernelπ_app_eq_zero}
  \leanok
  Let \(\iota : F \to G\) be a monomorphism of sheaves of \(R\)-modules on a site and let
  \(U\) be an object of the site. If \(c \in G(U)\) maps to zero in \((\coker \iota)(U)\),
  then \(c = \iota_U(a)\) for a (unique) \(a \in F(U)\).
\end{lemma}
\begin{proof}
  \leanok
  In an abelian category, a monomorphism is a kernel of its cokernel. Evaluation of sheaves
  at \(U\) preserves limits: it is the composite of the inclusion of sheaves into presheaves
  (a right adjoint, sheafification being left adjoint to it) with evaluation of presheaves at
  \(U\) (limits of presheaves are computed objectwise). Hence
  \(F(U) \to G(U)\) is a kernel of \(G(U) \to (\coker \iota)(U)\), which is the claim.
\end{proof}

\begin{lemma}[Monomorphisms are injective on sections]
  \label{lem:app_injective_of_mono}
  \lean{CategoryTheory.Sheaf.app_injective_of_mono}
  \leanok
  A monomorphism of sheaves of \(R\)-modules on a site is injective on sections over every
  object of the site.
\end{lemma}
\begin{proof}
  \leanok
  Evaluation of sheaves at an object preserves limits (as in
  \ref{lem:sections_left_exact}), hence preserves monomorphisms, and a monomorphism of
  \(R\)-modules is an injective map.
\end{proof}

\begin{theorem}[Sections surjectivity on an affine open]
  \label{thm:cokernel_app_surjective_affine_open}
  \lean{AlgebraicGeometry.IsAffineOpen.cokernel_app_surjective}
  \uses{def:moduleKSheaf}
  \leanok
  Let \(k\) be a commutative ring, \(X\) a scheme over \(\Spec k\), \(U\) an affine open of
  \(X\), \(G\) a sheaf of \(k\)-modules on the small Zariski site of \(X\), and
  \(\iota : \struct{X/k} \to G\) a monomorphism. Then every section of
  \(\coker \iota\) over \(U\) lifts to a section of \(G\) over \(U\): the map
  \(G(U) \to (\coker \iota)(U)\) is surjective.
\end{theorem}
\begin{proof}
  \leanok
  \uses{lem:sections_left_exact, lem:app_injective_of_mono, thm:affine_open_cech_cobounding}
  Let \(q \in (\coker \iota)(U)\) and write \(\pi : G \to \coker \iota\) for the projection.

  (a) \(\pi\) is an epimorphism of sheaves, hence locally surjective on sections: every point
  of \(U\) has an open neighbourhood on which the restriction of \(q\) lifts to \(G\).
  Shrinking, and since the basic opens of \(U\) form a basis of \(U\) (as \(U\) is affine),
  every point of \(U\) lies in some basic open \(D(g) \subseteq U\) carrying a lift of
  \(q|_{D(g)}\).

  (b) \(U\) is quasi-compact (it is affine), so there are finitely many sections
  \(f_1, \dots, f_n \in \Gamma(X, U)\) with \(D(f_1) \cup \dots \cup D(f_n) = U\) and lifts
  \(s_i \in G(D(f_i))\) with \(\pi(s_i) = q|_{D(f_i)}\).

  (c) On \(D(f_i) \cap D(f_j)\) the difference \(s_i - s_j\) is killed by \(\pi\), so by
  \ref{lem:sections_left_exact} there is a unique
  \(a_{ij} \in \Gamma\bigl(X, D(f_i) \cap D(f_j)\bigr)\) with
  \(\iota(a_{ij}) = s_i - s_j\).

  (d) The \(a_{ij}\) form a \v{C}ech \(1\)-cocycle: on triple intersections,
  \(\iota(a_{il}) = s_i - s_l = (s_i - s_j) + (s_j - s_l) = \iota(a_{ij} + a_{jl})\), and
  \(\iota\) is injective on sections (\ref{lem:app_injective_of_mono}).

  (e) By \ref{thm:affine_open_cech_cobounding} (applied to the affine open \(U\)) there are
  \(b_i \in \Gamma\bigl(X, D(f_i)\bigr)\) with \(a_{ij} = b_i - b_j\) on overlaps. The
  corrected lifts \(s_i - \iota(b_i)\) then agree on all overlaps, so by the sheaf condition
  they glue to a section \(s \in G(U)\).

  (f) Finally \(\pi(s) = q\): on each \(D(f_i)\) we have
  \(\pi(s) = \pi(s_i) - \pi(\iota(b_i)) = q|_{D(f_i)}\) since \(\pi \circ \iota = 0\), and
  sections of the sheaf \(\coker \iota\) that agree on a cover of \(U\) agree on \(U\).
\end{proof}

\begin{theorem}[Vanishing of \(H^1{}'\) on affine opens]
  \label{thm:affine_open_h1_vanishing}
  \lean{AlgebraicGeometry.IsAffineOpen.subsingleton_moduleKSheaf_hModule'_one}
  \uses{def:HModule', def:moduleKSheaf, thm:cokernel_app_surjective_affine_open}
  \leanok
  \source{stacks-project}
  Let \(k\) be a commutative ring, \(X\) a scheme over \(\Spec k\) and \(U\) an affine open
  of \(X\). Then
  \[
    H'^1\bigl(U, \struct{X/k}\bigr) \;=\; 0 .
  \]
  No hypothesis on \(X\) beyond the affineness of \(U\) is needed.
\end{theorem}
\begin{proof}
  \leanok
  \uses{lem:ext_one_vanishing_criterion, def:freeModuleSheafHomEquiv}
  The category of sheaves of \(k\)-modules on the small Zariski site of \(X\) is Grothendieck
  abelian, hence has enough injectives: choose a monomorphism
  \(\iota : \struct{X/k} \to G\) with \(G\) injective. By
  \ref{lem:ext_one_vanishing_criterion} (with \(L = k[U]\) the free sheaf of \(k\)-modules on
  \(U\)), it suffices to lift every morphism \(k[U] \to \coker \iota\) along
  \(G \to \coker \iota\). By \ref{def:freeModuleSheafHomEquiv}, morphisms out of \(k[U]\) are,
  compatibly with postcomposition, sections over \(U\); so the required lifting is exactly
  the surjectivity of \(G(U) \to (\coker \iota)(U)\), which is
  \ref{thm:cokernel_app_surjective_affine_open}.
\end{proof}

\begin{theorem}[Vanishing of \(H^1\) on affine schemes]
  \label{thm:affine_vanishing}
  \lean{AlgebraicGeometry.Scheme.subsingleton_moduleKSheaf_hModule_one}
  \uses{def:HModule, def:moduleKSheaf}
  \leanok
  \source{stacks-project}
  Let \(k\) be a commutative ring and \(X\) an affine scheme over \(\Spec k\). Then
  \[
    H^1\bigl(X, \struct{X/k}\bigr) \;=\; 0 .
  \]
\end{theorem}
\begin{proof}
  \leanok
  \uses{lem:ext_one_vanishing_criterion, thm:cokernel_app_surjective_affine_open,
    def:constHomEquiv, lem:constHomEquiv_naturality}
  The category of sheaves of \(k\)-modules on the small Zariski site of \(X\) is Grothendieck
  abelian, hence has enough injectives: choose a monomorphism
  \(\iota : \struct{X/k} \to G\) with \(G\) injective. By
  \ref{lem:ext_one_vanishing_criterion} (with \(L = k_J\) the constant sheaf), it suffices to
  lift every morphism \(k_J \to \coker \iota\) along \(G \to \coker \iota\). By
  \ref{def:constHomEquiv} and its naturality \ref{lem:constHomEquiv_naturality}, morphisms
  out of \(k_J\) are global sections, compatibly with postcomposition; so the required
  lifting is exactly the surjectivity of \(G(X) \to (\coker \iota)(X)\), i.e.\ of
  \(G(\top) \to (\coker \iota)(\top)\), which is
  \ref{thm:cokernel_app_surjective_affine_open} applied to the affine open \(U = \top\)
  (an affine scheme is an affine open of itself).
\end{proof}

\section{Twisted affine vanishing: quasi-coherent coefficients}
\label{sec:qcoh_vanishing}

The vanishing of the previous section is stated for the structure sheaf, but Serre's argument
only ever uses two features of \(\struct{X/k}\): that ambient sections act on sections, and that
sections localise on affine opens. This section isolates those two features as a packaging on an
arbitrary sheaf of \(k\)-modules and reruns the argument, obtaining degree-one vanishing on an
affine open with those coefficients. The intended application is to line bundles trivialised on
an affine chart, for which the packaging holds chart by chart.

For a sheaf \(F\) of \(k\)-modules on the small Zariski site of \(X\) and an inclusion of opens
\(W \subseteq V\), write \(s|_W\) (\emph{section restriction}) for the image of \(s \in F(V)\)
under the \(k\)-linear restriction \(F(V) \to F(W)\); it is functorial in the inclusion and
natural in \(F\).

\begin{definition}[Quasi-coherence packaging on an open]
  \label{def:qcohOn}
  \lean{AlgebraicGeometry.Scheme.QcohOn}
  \leanok
  Let \(X\) be a scheme, \(U\) an open of \(X\), and \(F\) a sheaf of \(k\)-modules on the small
  Zariski site of \(X\). A \emph{quasi-coherence packaging} of \(F\) on \(U\) is the data of:
  \begin{itemize}
    \item \textbf{(P1) an action} of the fixed ring \(\Gamma(X, U)\): for every open \(W \le U\)
      a scaling operation \(r \cdot m\) of \(r \in \Gamma(X, U)\) on \(m \in F(W)\), satisfying
      the module laws over \(\Gamma(X, U)\) (unit, associativity, and both distributivities) and
      natural under section restriction, \((r \cdot m)|_{W'} = r \cdot (m|_{W'})\) for \(W' \le
      W\);
    \item \textbf{(P2) localisation on affine opens} \(V \le U\), for every \(g \in \Gamma(X,
      U)\), in the two-lemma form of localisation away from \(g\):
      \begin{itemize}
        \item \emph{denominator clearing} — every \(c \in F(V \cap D(g))\) is \((g^N \cdot
          e)|_{V \cap D(g)}\) for some \(N \ge 0\) and \(e \in F(V)\);
        \item \emph{defect annihilation} — every \(t \in F(V)\) with \(t|_{V \cap D(g)} = 0\) is
          killed by a power of \(g\): \(g^M \cdot t = 0\) for some \(M \ge 0\).
      \end{itemize}
  \end{itemize}
  The action is by the single ring \(\Gamma(X, U)\) (not a per-open module structure), with the
  inclusion \(W \le U\) an explicit argument.
\end{definition}

\begin{theorem}[The structure sheaf is quasi-coherently packaged]
  \label{thm:moduleKSheaf_qcohOn}
  \lean{AlgebraicGeometry.Scheme.moduleKSheaf.instQcohOn}
  \uses{def:qcohOn, def:moduleKSheaf, lem:denominator_clearing, lem:section_annihilation}
  \leanok
  For \(X\) a scheme over \(\Spec k\), the structure sheaf \(\struct{X/k}\), viewed as a sheaf of
  \(k\)-modules, carries a quasi-coherence packaging on every open \(U\): the action of \(r \in
  \Gamma(X, U)\) on \(s \in \Gamma(X, W)\) is \(r|_W \cdot s\), the ordinary
  restriction-then-multiplication.
\end{theorem}
\begin{proof}
  \leanok
  The (P1) module laws are the ring axioms of \(\Gamma(X, W)\), and naturality under restriction
  holds because restriction is a ring homomorphism. The two (P2) axioms, with the action spelled
  out, are exactly denominator clearing (\ref{lem:denominator_clearing}) and annihilation
  (\ref{lem:section_annihilation}) for the structure sheaf.
\end{proof}

\begin{theorem}[Degree-one \v{C}ech cobounding, quasi-coherent coefficients]
  \label{thm:cech_cobounding_qcoh}
  \lean{AlgebraicGeometry.IsAffineOpen.exists_cech_cobounding_of_qcoh}
  \uses{def:qcohOn}
  \leanok
  Let \(U\) be an affine open of a scheme \(X\), covered by basic opens \(D(f_1), \dots, D(f_n)\),
  and let \(F\) be a sheaf of \(k\)-modules with a quasi-coherence packaging on \(U\). Every
  \v{C}ech \(1\)-cocycle \(a_{ij} \in F(D(f_i) \cap D(f_j))\) (with \(a_{il} = a_{ij} + a_{jl}\)
  on triple overlaps) is a coboundary: there are \(b_i \in F(D(f_i))\) with \(a_{ij} = b_i|_{D(f_i)
  \cap D(f_j)} - b_j|_{D(f_i) \cap D(f_j)}\).
\end{theorem}
\begin{proof}
  \leanok
  Verbatim the structure-sheaf argument \ref{thm:affine_open_cech_cobounding}, reading the
  packaging axioms in place of the structure-sheaf localisation. Clear denominators with one
  uniform exponent \(N\) (denominator clearing of the packaging), producing \(e_{ij} \in
  F(D(f_i))\) with \(e_{ij}|_{D(f_i) \cap D(f_j)} = f_j^N \cdot a_{ij}\). The triple-overlap
  defects \(t_{ijl} = e_{il} - e_{jl} - f_l^N \cdot a_{ij}\) vanish on the triple intersections by
  the cocycle identity, hence are killed by a uniform power \(f_l^M\) (defect annihilation). With
  a partition of unity \(1 = \sum_l h_l f_l^{N+M}\) of the affine \(U\), set \(b_i = \sum_l h_l
  f_l^M \cdot e_{il}\); the assembly \(b_i - b_j = \sum_l h_l f_l^M(e_{il} - e_{jl}) = (\sum_l h_l
  f_l^{N+M}) \cdot a_{ij} = a_{ij}\) uses only that the action distributes over restriction and
  sums, i.e.\ the (P1) mixin.
\end{proof}

\begin{theorem}[Sections surjectivity on an affine open, quasi-coherent coefficients]
  \label{thm:cokernel_app_surjective_qcoh}
  \lean{AlgebraicGeometry.IsAffineOpen.cokernel_app_surjective_of_qcoh}
  \uses{def:qcohOn, thm:cech_cobounding_qcoh, lem:sections_left_exact, lem:app_injective_of_mono}
  \leanok
  Let \(U\) be an affine open of a scheme \(X\), let \(F\) be a sheaf of \(k\)-modules with a
  quasi-coherence packaging on \(U\), and let \(\iota : F \to G\) be a monomorphism of sheaves of
  \(k\)-modules. Then every section of \(\coker \iota\) over \(U\) lifts to a section of \(G\)
  over \(U\).
\end{theorem}
\begin{proof}
  \leanok
  Verbatim \ref{thm:cokernel_app_surjective_affine_open} with \(F\) in place of \(\struct{X/k}\).
  Local lifts of \(q \in (\coker \iota)(U)\) exist on a finite basic cover \(D(f_i)\) of \(U\) by
  local surjectivity of the epimorphism \(\coker\iota\)-projection and quasi-compactness. On
  overlaps the differences of the lifts are killed by the projection, so by left-exactness of
  sections (\ref{lem:sections_left_exact}) they come from \(F\)-sections \(a_{ij}\); these form a
  \v{C}ech \(1\)-cocycle (\(\iota\) is injective on sections, \ref{lem:app_injective_of_mono}),
  which cobounds by \ref{thm:cech_cobounding_qcoh}. The corrected lifts glue to a section of \(G\)
  over \(U\) mapping to \(q\).
\end{proof}

\begin{theorem}[Twisted affine vanishing of \(H^1{}'\)]
  \label{thm:hModule1_vanishing_qcoh}
  \lean{AlgebraicGeometry.IsAffineOpen.subsingleton_hModule'_one_of_qcoh}
  \uses{def:qcohOn, def:HModule', thm:cokernel_app_surjective_qcoh,
    lem:ext_one_vanishing_criterion, def:freeModuleSheafHomEquiv}
  \leanok
  Let \(U\) be an affine open of a scheme \(X\) and \(F\) a sheaf of \(k\)-modules with a
  quasi-coherence packaging on \(U\). Then the degree-one cohomology of the object \(U\) with
  coefficients in \(F\) vanishes:
  \[
    H'^1(U, F) \;=\; 0 .
  \]
  Taking \(F = \struct{X/k}\) recovers \ref{thm:affine_open_h1_vanishing}.
\end{theorem}
\begin{proof}
  \leanok
  As in \ref{thm:affine_open_h1_vanishing}: choose a monomorphism \(\iota : F \to G\) with \(G\)
  injective; by the \(\Ext^1\) vanishing criterion \ref{lem:ext_one_vanishing_criterion} (with
  \(L = k[U]\) the free sheaf of \(k\)-modules on \(U\)) it suffices to lift every morphism
  \(k[U] \to \coker \iota\) along \(G \to \coker \iota\), which by \ref{def:freeModuleSheafHomEquiv}
  is the surjectivity of \(G(U) \to (\coker \iota)(U)\) — Theorem
  \ref{thm:cokernel_app_surjective_qcoh}.
\end{proof}

\begin{theorem}[Twisted affine vanishing of \(H^1\)]
  \label{thm:hModule_vanishing_qcoh}
  \lean{AlgebraicGeometry.Scheme.subsingleton_hModule_one_of_qcoh}
  \uses{def:qcohOn, def:HModule, thm:hModule1_vanishing_qcoh, def:HModule_linearEquivHModule'}
  \leanok
  Let \(X\) be an affine scheme and \(F\) a sheaf of \(k\)-modules with a quasi-coherence
  packaging on \(\top\). Then the degree-one cohomology of the site vanishes,
  \(H^1(X, F) = 0\). Taking \(F = \struct{X/k}\) recovers \ref{thm:affine_vanishing}.
\end{theorem}
\begin{proof}
  \leanok
  Apply \ref{thm:hModule1_vanishing_qcoh} to the affine open \(U = \top\) (an affine scheme is an
  affine open of itself) and transport the object-cohomology vanishing to the site cohomology
  along the equivalence \ref{def:HModule_linearEquivHModule'}.
\end{proof}

The two (P2) axioms of the packaging are not merely \emph{sufficient} for Serre's argument:
they say precisely that restriction from an affine chart to a basic open of the chart is a
localisation of modules. The remaining nodes of this section record that bridge, which lets
any packaged sheaf on an affine chart be read through the general theory of localised
modules --- in particular, it is how the two-lattice description of the two-cover complex
consumes the packaging.

\begin{definition}[The chart-ring module structure on packaged sections]
  \label{def:qcohOn_module}
  \lean{AlgebraicGeometry.Scheme.QcohOn.moduleOfLE}
  \uses{def:qcohOn}
  \leanok
  Let \(F\) carry a quasi-coherence packaging on \(U\) and let \(V \le U\). The (P1) action
  makes the \(k\)-module \(F(V)\) a module over the chart ring \(\Gamma(X, U)\): the unit,
  associativity and distributivity laws are the (P1) axioms, and the remaining law
  \(0 \cdot m = 0\) follows by adding \(0 \cdot m\) to both sides of
  \((0 + 0) \cdot m = 0 \cdot m\) and cancelling.
\end{definition}

\begin{definition}[Restriction is linear over the chart ring]
  \label{def:qcohOn_secRes_linear}
  \lean{AlgebraicGeometry.Scheme.QcohOn.secResₗ}
  \uses{def:qcohOn_module}
  \leanok
  For opens \(V \le W \le U\), section restriction \(F(W) \to F(V)\) is
  \(\Gamma(X, U)\)-linear for the module structures of \ref{def:qcohOn_module}: it is
  additive, being a map of \(k\)-modules, and it commutes with the action by the (P1)
  restriction-naturality \((r \cdot m)|_V = r \cdot (m|_V)\).
\end{definition}

\begin{theorem}[The packaging axioms are the localisation axioms]
  \label{thm:qcohOn_isLocalizedModule}
  \lean{AlgebraicGeometry.Scheme.QcohOn.isLocalizedModule_secResₗ}
  \uses{def:qcohOn, def:qcohOn_module, def:qcohOn_secRes_linear}
  \leanok
  Let \(U\) be an affine open of a scheme \(X\), let \(F\) be a sheaf of \(k\)-modules with
  a quasi-coherence packaging on \(U\), and let \(g \in \Gamma(X, U)\) act invertibly on
  \(F(U \cap D(g))\) (for the module structure of \ref{def:qcohOn_module}). Then the
  restriction map \(F(U) \to F(U \cap D(g))\) exhibits \(F(U \cap D(g))\) as the
  localisation of the \(\Gamma(X, U)\)-module \(F(U)\) at the multiplicative set
  \(S = \{1, g, g^2, \dots\}\) of powers of \(g\); that is:
  \begin{enumerate}
    \item every element of \(S\) acts invertibly on \(F(U \cap D(g))\);
    \item every \(c \in F(U \cap D(g))\) is of the form \(g^{-N} \cdot (e|_{U \cap D(g)})\)
      for some \(N \ge 0\) and \(e \in F(U)\);
    \item any two sections of \(F(U)\) with equal restriction to \(U \cap D(g)\) differ by
      an element killed by a power of \(g\).
  \end{enumerate}
  These three properties characterise the localisation map \(M \to S^{-1}M\) up to unique
  isomorphism under \(F(U)\).
\end{theorem}
\begin{proof}
  \leanok
  Each axiom is the term-for-term image of one piece of the packaging data. (1) Powers of
  an invertibly-acting element act invertibly. (2) Denominator clearing on the affine open
  \(U\) itself produces, for \(c \in F(U \cap D(g))\), an exponent \(N\) and a section
  \(e \in F(U)\) with \(e|_{U \cap D(g)} = g^N \cdot c\); acting by the inverse of
  \(g^N\) gives \(c = g^{-N} \cdot (e|_{U \cap D(g)})\). (3) If
  \(x_1|_{U \cap D(g)} = x_2|_{U \cap D(g)}\) then \(x_1 - x_2\) restricts to \(0\) on
  \(U \cap D(g)\), so defect annihilation yields \(M \ge 0\) with
  \(g^M \cdot (x_1 - x_2) = 0\).
\end{proof}

\begin{theorem}[Sections over a basic open are the localised sections]
  \label{thm:qcohOn_isLocalizedModule_moduleKSheaf}
  \lean{AlgebraicGeometry.Scheme.QcohOn.isLocalizedModule_secResₗ_moduleKSheaf}
  \uses{thm:qcohOn_isLocalizedModule, thm:moduleKSheaf_qcohOn, def:moduleKSheaf}
  \leanok
  Let \(X\) be a scheme over \(\Spec k\), \(U\) an affine open, and \(g \in \Gamma(X, U)\).
  Then restriction exhibits \(\Gamma(X, U \cap D(g))\) as the localisation of
  \(\Gamma(X, U)\) at the powers of \(g\) --- the classical
  \(\Gamma(X, D(g)) = \Gamma(X, U)_g\), with no hypothesis on \(g\).
\end{theorem}
\begin{proof}
  \leanok
  The structure sheaf carries the packaging on \(U\) (\ref{thm:moduleKSheaf_qcohOn}), with
  the action of \(r \in \Gamma(X, U)\) given by restriction-then-multiplication. The unit
  hypothesis of \ref{thm:qcohOn_isLocalizedModule} is discharged for free: \(g\) restricts
  to a unit on \(D(g)\), hence on \(U \cap D(g)\), so multiplication by
  \(g|_{U \cap D(g)}\) --- which is the packaging action of \(g\) --- is invertible on
  \(\Gamma(X, U \cap D(g))\). Apply \ref{thm:qcohOn_isLocalizedModule}.
\end{proof}

\section{The Mayer--Vietoris \((0,1)\)-slice}
\label{sec:mv-in-text}

For a Mayer--Vietoris square \(S\) (typically: opens \(X_1 = U \cap V\), \(X_2 = U\),
\(X_3 = V\), \(X_4 = U \cup V\) of a topological space) and a sheaf of \(R\)-modules \(F\) on
the site, this section produces the exactness of
\[
  0 \to F(X_4) \to F(X_2) \times F(X_3) \to F(X_1)
    \xrightarrow{\;\delta\;} H'^1(X_4, F) \to H'^1(X_2, F) \times H'^1(X_3, F)
\]
and derives the two-cover cokernel computation: if \(H'^1(X_2, F)\) and \(H'^1(X_3, F)\)
vanish, then \(F(X_1)/\mathrm{range}(\text{difference of restrictions}) \cong H'^1(X_4, F)\),
the quotient map being the connecting homomorphism \(\delta\).

\begin{definition}[Mayer--Vietoris square]
  \label{def:MayerVietorisSquare}
  \lean{CategoryTheory.GrothendieckTopology.MayerVietorisSquare}
  \mathlibok
  \textit{Provided by Mathlib.}
  Let \((\mathcal C, J)\) be a small site. A Mayer--Vietoris square consists of four objects
  \(X_1, X_2, X_3, X_4\) of \(\mathcal C\) and morphisms \(f_{12} : X_1 \to X_2\),
  \(f_{13} : X_1 \to X_3\), \(f_{24} : X_2 \to X_4\), \(f_{34} : X_3 \to X_4\) such that the
  two composites \(X_1 \to X_2 \to X_4\) and \(X_1 \to X_3 \to X_4\) agree, \(f_{13}\) is a
  monomorphism, and, after applying the representable-presheaf-then-sheafification functor
  \(\mathcal C \to \mathrm{Sh}(\mathcal C, J; \mathrm{Set})\), the resulting square of
  sheaves of sets is a pushout square: the image of \(X_4\) is the pushout, over the image of
  \(X_1\), of the images of \(X_2\) and \(X_3\). The prototypical example is
  \(X_1 = U \cap V\), \(X_2 = U\), \(X_3 = V\), \(X_4 = U \cup V\) for two opens \(U, V\) of a
  topological space, with the evident inclusions.
\end{definition}

\begin{definition}[The sheaf condition of a Mayer--Vietoris square]
  \label{def:mv_sheaf_condition}
  \lean{CategoryTheory.GrothendieckTopology.MayerVietorisSquare.SheafCondition,
    CategoryTheory.GrothendieckTopology.MayerVietorisSquare.sheafCondition_of_sheaf}
  \uses{def:MayerVietorisSquare}
  \mathlibok
  \textit{Provided by Mathlib.}
  Let \(S = (X_1,X_2,X_3,X_4,f_{12},f_{13},f_{24},f_{34})\) be a Mayer--Vietoris square on
  \((\mathcal C, J)\) and let \(P\) be a presheaf on \(\mathcal C\). Say \(P\) satisfies the
  sheaf condition for \(S\) when the square \(P(X_4) \to P(X_2) \to P(X_1)\),
  \(P(X_4) \to P(X_3) \to P(X_1)\) (restriction along \(f_{24}, f_{12}\) and along
  \(f_{34}, f_{13}\)) is a pullback square: equivalently, restriction identifies \(P(X_4)\)
  with the pairs \((u,v) \in P(X_2) \times P(X_3)\) of equal restriction to \(X_1\). Every
  sheaf satisfies the sheaf condition for every Mayer--Vietoris square of its site.
\end{definition}

\begin{definition}[The Mayer--Vietoris short complex of free sheaves]
  \label{def:moduleShortComplex}
  \lean{CategoryTheory.GrothendieckTopology.MayerVietorisSquare.moduleShortComplex}
  \uses{def:MayerVietorisSquare, def:freeModuleSheaf, def:freeModuleSheafMap}
  \leanok
  For a Mayer--Vietoris square \(S\) on \((\mathcal C, J)\), the complex of sheaves of
  \(R\)-modules
  \[
    R[X_1] \xrightarrow{\ (R[f_{12}],\, -R[f_{13}])\ } R[X_2] \oplus R[X_3]
      \xrightarrow{\ R[f_{24}] + R[f_{34}]\ } R[X_4],
  \]
  whose composite is zero by the commutativity of \(S\)'s square.
\end{definition}

\begin{theorem}[The Mayer--Vietoris short complex is short exact]
  \label{thm:moduleShortComplex_shortExact}
  \lean{CategoryTheory.GrothendieckTopology.MayerVietorisSquare.moduleShortComplex_shortExact}
  \uses{def:moduleShortComplex}
  \leanok
  For every Mayer--Vietoris square \(S\), the complex of \ref{def:moduleShortComplex} is a
  short exact sequence of sheaves of \(R\)-modules:
  \[
    0 \to R[X_1] \to R[X_2] \oplus R[X_3] \to R[X_4] \to 0 .
  \]
\end{theorem}
\begin{proof}
  \leanok
  \emph{The square of free sheaves is a pushout.} The free-\(R\)-module-sheaf functor
  \(U \mapsto R[U]\) (\ref{def:freeModuleSheafMap} on morphisms) is a left adjoint of the
  forgetful functor to sheaves of sets (it factors the free-\(R\)-module-on-a-set functor,
  itself a left adjoint, through sheafification, also a left adjoint). Left adjoints
  preserve colimits, in particular pushout squares; applying this functor to the defining
  pushout square of \(S\) (Definition~\ref{def:MayerVietorisSquare}: the square of
  sheafified representables at \(X_1,X_2,X_3,X_4\)) shows that the square
  \(R[X_1] \to R[X_2]\), \(R[X_1] \to R[X_3]\), \(R[X_2] \to R[X_4]\),
  \(R[X_3] \to R[X_4]\) is again a pushout, now in the category of sheaves of
  \(R\)-modules.

  \emph{Pushout along a mono is short exact.} In an additive category with biproducts, a
  commuting square with legs \(p : A \to B\), \(q : A \to C\), \(r : B \to D\),
  \(s : C \to D\) is a pushout square exactly when \(D\), together with the map
  \((r,s) : B \oplus C \to D\), is a cokernel of the difference map
  \((p,-q) : A \to B \oplus C\). Applied to the pushout square above, this shows that
  \(R[X_4]\), with the sum map, is a cokernel of the difference map
  \(R[X_1] \to R[X_2] \oplus R[X_3]\); that is, the complex of \ref{def:moduleShortComplex}
  is exact at \(R[X_2] \oplus R[X_3]\) and at \(R[X_4]\), with the map to \(R[X_4]\) an
  epimorphism.

  \emph{The difference map is a monomorphism.} By definition of a Mayer--Vietoris square,
  \(f_{13} : X_1 \to X_3\) is a monomorphism of \(\mathcal C\); postcomposition with a
  monomorphism is injective on morphism sets, so the representable-presheaf map it induces
  is a monomorphism of presheaves of sets, and both the free-\(R\)-module functor (on a set,
  mapping along an injection of index sets is injective on finitely supported functions) and
  sheafification preserve monomorphisms. Hence \(R[f_{13}] : R[X_1] \to R[X_3]\) is a
  monomorphism of sheaves of \(R\)-modules, and so is its negative. The second component of
  the difference map \((R[f_{12}], -R[f_{13}]) : R[X_1] \to R[X_2] \oplus R[X_3]\) (its
  composite with the projection to \(R[X_3]\)) is exactly \(-R[f_{13}]\); since a composite
  \(g \circ h\) being a monomorphism forces \(h\) to be a monomorphism, the difference map
  itself is a monomorphism.

  Together, the three paragraphs give short exactness.
\end{proof}

\begin{definition}[The restriction-difference map]
  \label{def:moduleDiff}
  \lean{CategoryTheory.GrothendieckTopology.MayerVietorisSquare.moduleDiff}
  \uses{def:MayerVietorisSquare}
  \leanok
  For a Mayer--Vietoris square \(S\) and a sheaf \(F\) of \(R\)-modules, the \(R\)-linear map
  \[
    F(X_2) \times F(X_3) \longrightarrow F(X_1), \qquad
    (s_2, s_3) \longmapsto s_2|_{X_1} - s_3|_{X_1},
  \]
  restricting along \(f_{12}\) and \(f_{13}\) respectively and taking the difference.
\end{definition}

\begin{lemma}[Degree-zero exactness]
  \label{lem:mv_exact_degree0}
  \lean{CategoryTheory.GrothendieckTopology.MayerVietorisSquare.sections_ext,
    CategoryTheory.GrothendieckTopology.MayerVietorisSquare.exists_glue_of_moduleDiff_eq_zero}
  \uses{def:moduleDiff, def:mv_sheaf_condition}
  \leanok
  For a Mayer--Vietoris square \(S\) and a sheaf \(F\) of \(R\)-modules, the sequence
  \[
    0 \to F(X_4) \xrightarrow{\ (\mathrm{res},\,\mathrm{res})\ } F(X_2) \times F(X_3)
      \xrightarrow{\ \mathrm{moduleDiff}\ } F(X_1)
  \]
  is exact: the pair of restrictions to \(X_2, X_3\) is injective, and its image is exactly
  the kernel of \ref{def:moduleDiff}.
\end{lemma}
\begin{proof}
  \leanok
  This is the sheaf condition of \(S\) for (the sheaf of sets underlying) \(F\)
  (Definition~\ref{def:mv_sheaf_condition}): the map \(F(X_4) \to F(X_2) \times F(X_3)\) has
  image exactly the pairs with equal restriction to \(X_1\) (i.e.\
  \(\ker(\mathrm{moduleDiff})\)), and two elements of \(F(X_4)\) with the same pair of
  restrictions are equal.
\end{proof}

\begin{lemma}[Contravariant long exact sequence of \(\Ext\)]
  \label{lem:ext_contravariant_les}
  \lean{CategoryTheory.Abelian.Ext.contravariant_sequence_exact₁,
    CategoryTheory.Abelian.Ext.contravariant_sequence_exact₃,
    CategoryTheory.ShortComplex.ShortExact.extClass,
    CategoryTheory.Abelian.Ext.biprodAddEquiv}
  \mathlibok
  \textit{Provided by Mathlib.}
  Let \(0 \to A \xrightarrow{f} B \xrightarrow{g} C \to 0\) be a short exact sequence in an
  abelian category with \(\Ext\)-groups and let \(Y\) be any object. There is a canonical
  extension class \(\delta \in \Ext^1(C, A)\) (built from the connecting homomorphism of the
  sequence), and for every \(n \ge 0\) the sequence of abelian groups
  \[
    \Ext^n(C,Y) \xrightarrow{g^*} \Ext^n(B,Y) \xrightarrow{f^*} \Ext^n(A,Y)
      \xrightarrow{\ \delta \cdot (-)\ } \Ext^{n+1}(C,Y) \xrightarrow{g^*} \Ext^{n+1}(B,Y)
      \xrightarrow{f^*} \Ext^{n+1}(A,Y)
  \]
  is exact at each of the four inner terms, where \(f^*, g^*\) are precomposition
  (contravariant functoriality of \(\Ext(-,Y)\)) and \(\delta \cdot (-)\) is composition
  (Yoneda product) with \(\delta\). (Convention: \(\Ext^0 = \Hom\).) If \(B = X_2 \oplus X_3\)
  is a direct sum, \(\Ext^n(B,Y) \cong \Ext^n(X_2,Y) \times \Ext^n(X_3,Y)\) naturally,
  compatibly with \(f^*, g^*\).
\end{lemma}

\begin{definition}[The connecting map]
  \label{def:moduleDelta}
  \lean{CategoryTheory.GrothendieckTopology.MayerVietorisSquare.moduleDelta}
  \uses{def:HModule'_linearEquiv0, thm:moduleShortComplex_shortExact,
    lem:ext_contravariant_les}
  \leanok
  For a Mayer--Vietoris square \(S\) and a sheaf \(F\) of \(R\)-modules, the connecting map
  \(F(X_1) \to H'^1(X_4, F)\): the identification
  \(F(X_1) \cong H'^0(X_1,F) = \Ext^0(R[X_1], F)\) of \ref{def:HModule'_linearEquiv0},
  followed by the map \(\delta \cdot (-) : \Ext^0(R[X_1],F) \to \Ext^1(R[X_4],F) = H'^1(X_4,F)\)
  of \ref{lem:ext_contravariant_les} for the short exact sequence of
  \ref{thm:moduleShortComplex_shortExact} (\(A = R[X_1]\), \(B = R[X_2] \oplus R[X_3]\),
  \(C = R[X_4]\)).
\end{definition}

\begin{lemma}[Exactness at \(F(X_1)\)]
  \label{lem:mv_exact_at_X1}
  \lean{CategoryTheory.GrothendieckTopology.MayerVietorisSquare.moduleDelta_moduleDiff,
    CategoryTheory.GrothendieckTopology.MayerVietorisSquare.exists_of_moduleDelta_eq_zero}
  \uses{def:moduleDelta, def:moduleDiff}
  \leanok
  For a Mayer--Vietoris square \(S\) and a sheaf \(F\) of \(R\)-modules,
  \(\ker(\mathrm{moduleDelta}) = \mathrm{range}(\mathrm{moduleDiff})\) as submodules of
  \(F(X_1)\): the connecting class of a difference of restrictions vanishes, and conversely
  every section of \(F(X_1)\) with vanishing connecting class is such a difference.
\end{lemma}
\begin{proof}
  \leanok
  \uses{lem:ext_contravariant_les, thm:moduleShortComplex_shortExact,
    lem:freeModuleSheafHomEquiv_naturality}
  This is exactness of the sequence of \ref{lem:ext_contravariant_les} (for \(n=0\), the
  short exact sequence of \ref{thm:moduleShortComplex_shortExact}, \(Y=F\)) at the term
  \(\Ext^0(A,F) = F(X_1)\), transported along the identifications
  \(\Ext^0(B,F) \cong F(X_2) \times F(X_3)\) (\ref{lem:ext_contravariant_les}, direct sum) and
  \(\Ext^0(A,F) \cong F(X_1)\) (\ref{def:HModule'_linearEquiv0}): under these identifications
  \(f^* : \Ext^0(B,F) \to \Ext^0(A,F)\) becomes \(\mathrm{moduleDiff}\) and
  \(\delta \cdot (-) : \Ext^0(A,F) \to \Ext^1(C,F)\) becomes \(\mathrm{moduleDelta}\), by
  naturality of the section equivalence in the site variable
  (\ref{lem:freeModuleSheafHomEquiv_naturality}), which identifies precomposition at the
  level of morphisms out of a free sheaf with restriction at the level of sections.
\end{proof}

\begin{lemma}[Exactness at \(H'^1(X_4,F)\)]
  \label{lem:mv_exact_at_H1X4}
  \lean{CategoryTheory.GrothendieckTopology.MayerVietorisSquare.res_moduleDelta₂,
    CategoryTheory.GrothendieckTopology.MayerVietorisSquare.res_moduleDelta₃,
    CategoryTheory.GrothendieckTopology.MayerVietorisSquare.exists_moduleDelta_eq}
  \uses{def:moduleDelta, def:HModule'_res}
  \leanok
  For a Mayer--Vietoris square \(S\) and a sheaf \(F\) of \(R\)-modules,
  \(\ker(\mathrm{res}_{f_{24}}) \cap \ker(\mathrm{res}_{f_{34}}) = \mathrm{range}(\mathrm{moduleDelta})\)
  as submodules of \(H'^1(X_4,F)\), where
  \(\mathrm{res}_{f_{24}}, \mathrm{res}_{f_{34}} : H'^1(X_4,F) \to H'^1(X_2,F), H'^1(X_3,F)\)
  are the restrictions of \ref{def:HModule'_res}: the connecting class of any section of
  \(F(X_1)\) restricts to zero on both \(X_2\) and \(X_3\), and conversely every class
  restricting to zero on both is a connecting class.
\end{lemma}
\begin{proof}
  \leanok
  \uses{lem:ext_contravariant_les, thm:moduleShortComplex_shortExact}
  As in \ref{lem:mv_exact_at_X1}, this is exactness of the sequence of
  \ref{lem:ext_contravariant_les} (for \(n=0\), i.e.\ at the term \(\Ext^1(C,F)\) between
  \(\delta \cdot (-)\) and \(g^*\)) transported along the same identifications, together
  with the splitting \(\Ext^1(B,F) \cong H'^1(X_2,F) \times H'^1(X_3,F)\) under which
  \(g^*\) becomes the pair \((\mathrm{res}_{f_{24}}, \mathrm{res}_{f_{34}})\).
\end{proof}

\begin{theorem}[Surjectivity of the connecting map]
  \label{thm:moduleDelta_surjective}
  \lean{CategoryTheory.GrothendieckTopology.MayerVietorisSquare.moduleDelta_surjective}
  \uses{lem:mv_exact_at_H1X4}
  \leanok
  If \(H'^1(X_2,F)\) and \(H'^1(X_3,F)\) are both zero, then
  \(\mathrm{moduleDelta} : F(X_1) \to H'^1(X_4,F)\) is surjective.
\end{theorem}
\begin{proof}
  \leanok
  Immediate from \ref{lem:mv_exact_at_H1X4}:
  \(\ker(\mathrm{res}_{f_{24}}) \cap \ker(\mathrm{res}_{f_{34}}) = H'^1(X_4,F)\) once both
  targets vanish, so \(\mathrm{range}(\mathrm{moduleDelta}) = H'^1(X_4,F)\).
\end{proof}

\begin{definition}[The connecting map on the cokernel]
  \label{def:moduleDeltaQuotient}
  \lean{CategoryTheory.GrothendieckTopology.MayerVietorisSquare.moduleDeltaQuotient}
  \uses{def:moduleDelta, lem:mv_exact_at_X1}
  \leanok
  Since \(\mathrm{range}(\mathrm{moduleDiff}) \subseteq \ker(\mathrm{moduleDelta})\)
  (\ref{lem:mv_exact_at_X1}), \(\mathrm{moduleDelta}\) descends to an \(R\)-linear map
  \[
    F(X_1)/\mathrm{range}(\mathrm{moduleDiff}) \longrightarrow H'^1(X_4,F).
  \]
\end{definition}

\begin{theorem}[The Mayer--Vietoris two-cover cokernel computation]
  \label{thm:h1LinearEquiv}
  \lean{CategoryTheory.GrothendieckTopology.MayerVietorisSquare.h1LinearEquiv}
  \uses{def:moduleDeltaQuotient, lem:mv_exact_at_X1, thm:moduleDelta_surjective}
  \leanok
  \source{stacks-project}
  If \(H'^1(X_2,F)\) and \(H'^1(X_3,F)\) are both zero, the map of
  \ref{def:moduleDeltaQuotient} is an \(R\)-linear isomorphism
  \[
    F(X_1)/\mathrm{range}(\mathrm{moduleDiff}) \;\cong\; H'^1(X_4,F).
  \]
\end{theorem}
\begin{proof}
  \leanok
  Injectivity: if the class of \(s \in F(X_1)\) maps to \(0\), then
  \(\mathrm{moduleDelta}(s) = 0\), so by \ref{lem:mv_exact_at_X1}
  \(s \in \mathrm{range}(\mathrm{moduleDiff})\), i.e.\ the class of \(s\) is already \(0\).
  Surjectivity is \ref{thm:moduleDelta_surjective}.
\end{proof}

\begin{lemma}[Compatibility with the connecting map]
  \label{lem:h1LinearEquiv_mk}
  \lean{CategoryTheory.GrothendieckTopology.MayerVietorisSquare.h1LinearEquiv_mk}
  \uses{thm:h1LinearEquiv}
  \leanok
  Under the isomorphism of \ref{thm:h1LinearEquiv}, the class of \(s \in F(X_1)\)
  corresponds to \(\mathrm{moduleDelta}(s)\).
\end{lemma}
\begin{proof}
  \leanok
  Immediate from the construction of \ref{def:moduleDeltaQuotient} and
  \ref{thm:h1LinearEquiv}.
\end{proof}

\section{The two-affine-cover cokernel bridge}

For a scheme \(X\) over \(\Spec k\) covered by two affine opens \(U_0 \cup U_1 = X\), this
section computes the degree-one cohomology of the structure sheaf as
\[
  H^1\bigl(X, \struct{X/k}\bigr) \;\cong\;
    \Gamma(X, U_0 \cap U_1) \big/
      \mathrm{range}\bigl((s_0,s_1) \mapsto s_0|_{U_0 \cap U_1} - s_1|_{U_0 \cap U_1}\bigr),
\]
the computational access point for the genus of a curve. No hypothesis on the intersection
\(U_0 \cap U_1\) is needed: the derived Mayer--Vietoris sequence of Section~\ref{sec:mv-in-text}
is exact regardless, and affineness of the two pieces alone kills their degree-one
cohomology (\ref{thm:affine_open_h1_vanishing}) --- the classical affine-intersection
hypothesis is needed only for the \v{C}ech-complex identification of \(H^1\), which this
route bypasses.

\begin{definition}[Mayer--Vietoris square of two opens]
  \label{def:opens_mayerVietorisSquare'}
  \lean{Opens.mayerVietorisSquare'}
  \uses{def:MayerVietorisSquare}
  \mathlibok
  \textit{Provided by Mathlib.}
  For a topological space \(T\) and opens \(X_1,X_2,X_3,X_4 \subseteq T\) with
  \(X_1 = X_2 \cap X_3\) and \(X_4 = X_2 \cup X_3\), the square \((X_1,X_2,X_3,X_4)\) with the
  inclusion morphisms is a Mayer--Vietoris square (Definition~\ref{def:MayerVietorisSquare})
  on the small Zariski site of \(T\).
\end{definition}

\begin{definition}[The two-cover square]
  \label{def:twoCoverSquare}
  \lean{AlgebraicGeometry.Scheme.twoCoverSquare}
  \uses{def:opens_mayerVietorisSquare'}
  \leanok
  For a scheme \(X\) and opens \(U_0, U_1 \subseteq X\) with \(U_0 \cup U_1 = X\), the
  Mayer--Vietoris square \((U_0 \cap U_1,\, U_0,\, U_1,\, X)\) on the small Zariski site of
  \(X\) (Definition~\ref{def:opens_mayerVietorisSquare'}, specialized to \(X_4 = X\)), with
  the inclusion morphisms.
\end{definition}

\begin{theorem}[The two-cover \(H^1\) computation, general coefficients]
  \label{thm:twoCoverH1LinearEquiv}
  \lean{AlgebraicGeometry.Scheme.twoCoverH1LinearEquiv}
  \uses{def:twoCoverSquare, def:HModule_linearEquivHModule', thm:h1LinearEquiv}
  \leanok
  \source{stacks-project}
  Let \(k\) be a commutative ring, \(X\) a scheme over \(\Spec k\), \(U_0, U_1 \subseteq X\)
  opens with \(U_0 \cup U_1 = X\), and \(F\) a sheaf of \(k\)-modules on the small Zariski
  site of \(X\) with \(H'^1(U_0,F) = H'^1(U_1,F) = 0\). Then
  \[
    H^1(X,F) \;\cong\;
      F(U_0 \cap U_1) \big/ \mathrm{range}\bigl((s_0,s_1) \mapsto
        s_0|_{U_0 \cap U_1} - s_1|_{U_0 \cap U_1}\bigr) .
  \]
\end{theorem}
\begin{proof}
  \leanok
  Compose the isomorphism \(H^1(X,F) \cong H'^1(X,F)\) of
  \ref{def:HModule_linearEquivHModule'} (\(X\) is a terminal object of its own small Zariski
  site) with the inverse of the isomorphism of \ref{thm:h1LinearEquiv} for the two-cover
  square of \ref{def:twoCoverSquare}, whose \(\mathrm{moduleDiff}\) is exactly the displayed
  restriction-difference map.
\end{proof}

\begin{definition}[The two-cover restriction-difference map]
  \label{def:TwoCover_diff}
  \lean{AlgebraicGeometry.TwoCover.diff}
  \uses{def:moduleKSheaf}
  \leanok
  Let \(k\) be a commutative ring, \(X\) a scheme over \(\Spec k\), and \(U_0, U_1 \subseteq X\)
  opens. The \(k\)-linear map
  \[
    \Gamma(X,U_0) \times \Gamma(X,U_1) \longrightarrow \Gamma(X, U_0 \cap U_1), \qquad
    (s_0,s_1) \longmapsto s_0|_{U_0 \cap U_1} - s_1|_{U_0 \cap U_1} .
  \]
  When \(U_0 \cup U_1 = X\), this is the map \ref{def:moduleDiff} of the two-cover square
  (Definition~\ref{def:twoCoverSquare}) for the structure sheaf \(\struct{X/k}\).
\end{definition}

\begin{definition}[The two-cover cokernel carrier]
  \label{def:TwoCover_H1Cok}
  \lean{AlgebraicGeometry.TwoCover.H1Cok}
  \uses{def:TwoCover_diff}
  \leanok
  The \(k\)-module \(\Gamma(X, U_0 \cap U_1)/\mathrm{range}(\mathrm{diff})\), the cokernel of
  the map of \ref{def:TwoCover_diff}.
\end{definition}

\begin{lemma}[Sections extending to one piece vanish in the cokernel]
  \label{lem:TwoCover_h1Cok_mk_res}
  \lean{AlgebraicGeometry.TwoCover.h1Cok_mk_resHom_left,
    AlgebraicGeometry.TwoCover.h1Cok_mk_resHom_right}
  \uses{def:TwoCover_H1Cok}
  \leanok
  For \(t \in \Gamma(X,U_0)\), the class of \(t|_{U_0 \cap U_1}\) in
  \ref{def:TwoCover_H1Cok} is zero; symmetrically, for \(t \in \Gamma(X,U_1)\) the class of
  \(t|_{U_0 \cap U_1}\) is zero.
\end{lemma}
\begin{proof}
  \leanok
  The first is the class of \(\mathrm{diff}(t,0)\), the second of \(\mathrm{diff}(0,-t)\)
  (Definition~\ref{def:TwoCover_diff}), both in the range of \(\mathrm{diff}\).
\end{proof}

\begin{definition}[The two-cover connecting map]
  \label{def:TwoCover_delta}
  \lean{AlgebraicGeometry.TwoCover.delta}
  \uses{def:twoCoverSquare, def:HModule_linearEquivHModule', def:moduleDelta}
  \leanok
  Assume \(U_0 \cup U_1 = X\). The \(k\)-linear map
  \(\Gamma(X, U_0 \cap U_1) \to H^1(X, \struct{X/k})\): the connecting map
  \ref{def:moduleDelta} of the two-cover square (Definition~\ref{def:twoCoverSquare}) for
  \(F = \struct{X/k}\), followed by the inverse of the isomorphism
  \(H^1(X,\struct{X/k}) \cong H'^1(X,\struct{X/k})\) of
  \ref{def:HModule_linearEquivHModule'}.
\end{definition}

\begin{lemma}[The connecting class of a difference vanishes]
  \label{lem:TwoCover_delta_diff}
  \lean{AlgebraicGeometry.TwoCover.delta_diff}
  \uses{def:TwoCover_delta, def:TwoCover_diff, lem:mv_exact_at_X1}
  \leanok
  For \(t \in \Gamma(X,U_0) \times \Gamma(X,U_1)\), \(\mathrm{delta}(\mathrm{diff}(t)) = 0\).
\end{lemma}
\begin{proof}
  \leanok
  \(\mathrm{diff}(t)\) equals the two-cover square's \(\mathrm{moduleDiff}(t)\)
  (Definition~\ref{def:TwoCover_diff}), whose connecting class vanishes by
  \ref{lem:mv_exact_at_X1}; transport along the (linear, hence class-zero-preserving)
  isomorphism of \ref{def:HModule_linearEquivHModule'} used in the definition of
  \(\mathrm{delta}\).
\end{proof}

\begin{lemma}[The connecting class of an extending section vanishes]
  \label{lem:TwoCover_delta_resHom}
  \lean{AlgebraicGeometry.TwoCover.delta_resHom_left,
    AlgebraicGeometry.TwoCover.delta_resHom_right}
  \uses{lem:TwoCover_delta_diff}
  \leanok
  For \(t \in \Gamma(X,U_0)\), \(\mathrm{delta}(t|_{U_0 \cap U_1}) = 0\); symmetrically, for
  \(t \in \Gamma(X,U_1)\), \(\mathrm{delta}(t|_{U_0 \cap U_1}) = 0\).
\end{lemma}
\begin{proof}
  \leanok
  \(t|_{U_0 \cap U_1} = \mathrm{diff}(t,0)\), resp.\ \(-\mathrm{diff}(0,t)\)
  (Definition~\ref{def:TwoCover_diff}), so this is \ref{lem:TwoCover_delta_diff} (using
  \(\mathrm{delta}\)'s linearity in the second case).
\end{proof}

\begin{theorem}[Surjectivity of the two-cover connecting map]
  \label{thm:TwoCover_delta_surjective}
  \lean{AlgebraicGeometry.TwoCover.delta_surjective}
  \uses{def:TwoCover_delta, thm:affine_open_h1_vanishing, thm:moduleDelta_surjective}
  \leanok
  If \(U_0\) and \(U_1\) are affine opens of \(X\) (still with \(U_0 \cup U_1 = X\)), then
  \(\mathrm{delta} : \Gamma(X, U_0 \cap U_1) \to H^1(X,\struct{X/k})\) is surjective.
\end{theorem}
\begin{proof}
  \leanok
  By \ref{thm:affine_open_h1_vanishing}, \(H'^1(U_0,\struct{X/k}) = H'^1(U_1,\struct{X/k}) = 0\);
  by \ref{thm:moduleDelta_surjective} applied to the two-cover square, its connecting map
  \(\mathrm{moduleDelta}\) is surjective onto \(H'^1(X,\struct{X/k})\), and transporting
  along the isomorphism of \ref{def:HModule_linearEquivHModule'} used to define
  \(\mathrm{delta}\) preserves surjectivity.
\end{proof}

\begin{theorem}[The two-affine-cover \(H^1\) cokernel bridge]
  \label{thm:TwoCover_h1CokEquiv}
  \lean{AlgebraicGeometry.TwoCover.h1CokEquiv}
  \uses{thm:affine_open_h1_vanishing, thm:h1LinearEquiv, def:HModule_linearEquivHModule',
    def:TwoCover_H1Cok}
  \leanok
  \source{stacks-project}
  Let \(k\) be a commutative ring, \(X\) a scheme over \(\Spec k\), and \(U_0, U_1\) affine
  opens of \(X\) with \(U_0 \cup U_1 = X\). There is a \(k\)-linear isomorphism
  \[
    H^1\bigl(X, \struct{X/k}\bigr) \;\cong\; \Gamma(X, U_0 \cap U_1)\big/\mathrm{range}(\mathrm{diff}),
  \]
  where \(\mathrm{diff}(s_0,s_1) = s_0|_{U_0 \cap U_1} - s_1|_{U_0 \cap U_1}\)
  (\ref{def:TwoCover_diff}). No hypothesis on \(U_0 \cap U_1\) (affineness, separatedness, or
  even quasi-compactness) is needed: the derived Mayer--Vietoris sequence of
  \ref{thm:moduleShortComplex_shortExact} is exact regardless, and affineness of \(U_0, U_1\)
  alone kills their degree-one cohomology (\ref{thm:affine_open_h1_vanishing}); the classical
  hypothesis on the intersection is needed only for the \v{C}ech-complex identification of
  \(H^1\), which this route bypasses.
\end{theorem}
\begin{proof}
  \leanok
  By \ref{thm:affine_open_h1_vanishing}, \(H'^1(U_0,\struct{X/k}) = H'^1(U_1,\struct{X/k}) = 0\).
  Apply \ref{thm:h1LinearEquiv} to the two-cover square of \ref{def:twoCoverSquare} with
  \(F = \struct{X/k}\), giving
  \(\Gamma(X, U_0 \cap U_1)/\mathrm{range}(\mathrm{diff}) \cong H'^1(X,\struct{X/k})\), and
  compose with the inverse of the isomorphism
  \(H^1(X,\struct{X/k}) \cong H'^1(X,\struct{X/k})\) of
  \ref{def:HModule_linearEquivHModule'}.
\end{proof}

\begin{lemma}[Compatibility with the connecting map]
  \label{lem:TwoCover_h1CokEquiv_delta}
  \lean{AlgebraicGeometry.TwoCover.h1CokEquiv_delta}
  \uses{thm:TwoCover_h1CokEquiv, def:TwoCover_delta, lem:h1LinearEquiv_mk}
  \leanok
  Under the isomorphism of \ref{thm:TwoCover_h1CokEquiv}, the class
  \(\mathrm{delta}(s)\) (Definition~\ref{def:TwoCover_delta}) corresponds, for every
  \(s \in \Gamma(X, U_0 \cap U_1)\), to the class of \(s\) in
  \(\Gamma(X,U_0 \cap U_1)/\mathrm{range}(\mathrm{diff})\).
\end{lemma}
\begin{proof}
  \leanok
  Unwind the definitions of \ref{thm:TwoCover_h1CokEquiv} and \ref{def:TwoCover_delta} (both
  built from the same two ingredients, \ref{thm:h1LinearEquiv} and
  \ref{def:HModule_linearEquivHModule'}) and apply the compatibility \ref{lem:h1LinearEquiv_mk}
  of the Mayer--Vietoris cokernel isomorphism with its connecting map.
\end{proof}

\begin{lemma}[Compatibility with the connecting map, inverse form]
  \label{lem:TwoCover_h1CokEquiv_symm_mk}
  \lean{AlgebraicGeometry.TwoCover.h1CokEquiv_symm_mk}
  \uses{lem:TwoCover_h1CokEquiv_delta}
  \leanok
  The inverse of the isomorphism of \ref{thm:TwoCover_h1CokEquiv} sends the class of
  \(s \in \Gamma(X, U_0 \cap U_1)\) to \(\mathrm{delta}(s)\).
\end{lemma}
\begin{proof}
  \leanok
  Apply the isomorphism of \ref{thm:TwoCover_h1CokEquiv} to both sides of
  \ref{lem:TwoCover_h1CokEquiv_delta} and use injectivity.
\end{proof}

\section{Base change of the two-cover cohomology over an affine test ring}

Fix a field \(k\) and an object \(C\) of \(\Over(\Spec k)\). For a commutative \(k\)-algebra \(R\)
--- the \emph{affine test ring} --- base change produces the relative curve
\(C_R = C \otimes \Spec R\) over \(\Spec R\); an affine two-chart cover of \(C\)
(Section~\ref{sec:mv-in-text} and Chapter~\ref{chap:Curves}) pulls back to one of \(C_R\), and the
two-affine-cover bridge of the previous section then computes \(H^1(C_R, \struct{C_R/R})\) as an
explicit two-cover cokernel of \(R\)-modules. When \(C\) is a smooth curve, the degree-zero
cohomology of the relative curve is the test ring itself, \(H^0(C_R, \struct{C_R/R}) = R\).

\begin{definition}[The relative curve over the test ring]
  \label{def:relCurve}
  \lean{AlgebraicGeometry.relCurve, AlgebraicGeometry.relCurve.instOver}
  \leanok
  For \(C \in \Over(\Spec k)\) and a commutative \(k\)-algebra \(R\), the \emph{relative curve}
  \(C_R\) is the underlying scheme of the fibre product \(C \otimes \Spec R\) of \(C\) with the
  object \(\Spec R\) of \(\Over(\Spec k)\), regarded as a scheme over \(\Spec R\) via the second
  projection \(\mathrm{pr}_2 : C \otimes \Spec R \to \Spec R\).
\end{definition}

\begin{definition}[The base-changed affine two-cover]
  \label{def:relCover}
  \lean{AlgebraicGeometry.relCover}
  \uses{def:relCurve, def:AffineTwoCover, thm:AffineTwoCover_pullbackProd}
  \leanok
  Let \(D = (V_0, V_1)\) be an affine two-chart cover (\ref{def:AffineTwoCover}) of the underlying
  scheme of \(C\). The \emph{base-changed cover} of the relative curve \(C_R\) is the pair of
  preimages \(\bigl(\mathrm{pr}_1^{-1} V_0,\ \mathrm{pr}_1^{-1} V_1\bigr)\) of the charts of \(D\)
  under the first projection \(\mathrm{pr}_1 : C_R \to C\); by
  \ref{thm:AffineTwoCover_pullbackProd} it is again an affine two-chart cover, now of \(C_R\).
\end{definition}

\begin{lemma}[The base-changed charts are affine and cover the relative curve]
  \label{lem:relCover_charts}
  \lean{AlgebraicGeometry.relCover_isAffineOpen₀, AlgebraicGeometry.relCover_isAffineOpen₁,
    AlgebraicGeometry.relCover_sup}
  \uses{def:relCover}
  \leanok
  The two charts \(\mathrm{pr}_1^{-1} V_0\) and \(\mathrm{pr}_1^{-1} V_1\) of the base-changed cover
  are affine open, and their union is all of \(C_R\).
\end{lemma}
\begin{proof}
  \leanok
  \uses{thm:AffineTwoCover_pullbackProd}
  These are the affineness of the two charts and the covering property recorded in the
  affine-two-chart-cover structure \(\mathrm{pr}_1^{-1} D\) of \ref{thm:AffineTwoCover_pullbackProd}.
\end{proof}

\begin{lemma}[The overlap of the base-changed charts]
  \label{lem:relCover_inf}
  \lean{AlgebraicGeometry.relCover_inf}
  \uses{def:relCover}
  \leanok
  The overlap of the two base-changed charts is the preimage of the overlap of \(D\):
  \[
    \mathrm{pr}_1^{-1} V_0 \cap \mathrm{pr}_1^{-1} V_1 = \mathrm{pr}_1^{-1}(V_0 \cap V_1).
  \]
\end{lemma}
\begin{proof}
  \leanok
  Taking preimages under \(\mathrm{pr}_1\) commutes with intersection of opens.
\end{proof}

\begin{definition}[The relative two-cover cohomology carrier]
  \label{def:relTwoCoverH1}
  \lean{AlgebraicGeometry.relTwoCoverH1}
  \uses{def:relCurve, def:relCover, def:moduleKSheaf, def:TwoCover_diff, def:TwoCover_H1Cok}
  \leanok
  For the structure sheaf of the relative curve \(C_R\) viewed as a sheaf of \(R\)-modules, its
  degree-one cohomology over \(R\) is the two-cover cokernel of the base-changed cover: an
  \(R\)-linear isomorphism
  \[
    H^1\bigl(C_R, \struct{C_R/R}\bigr) \;\cong\;
      \Gamma\bigl(C_R,\ \mathrm{pr}_1^{-1}V_0 \cap \mathrm{pr}_1^{-1}V_1\bigr)\big/
        \mathrm{range}(\mathrm{diff}),
  \]
  where \(\mathrm{diff}(s_0,s_1) = s_0|_\cap - s_1|_\cap\) (\ref{def:TwoCover_diff}).
\end{definition}
\begin{proof}
  \leanok
  \uses{thm:TwoCover_h1CokEquiv, lem:relCover_charts}
  The base-changed charts are affine and cover \(C_R\) (\ref{lem:relCover_charts}), so this is the
  two-affine-cover \(H^1\) cokernel bridge \ref{thm:TwoCover_h1CokEquiv} applied over the base ring
  \(R\) to the scheme \(C_R\) with its base-changed cover; the affine-vanishing inputs required by
  that bridge are exactly the affineness of the two charts.
\end{proof}

\begin{theorem}[Degree-zero cohomology of the relative curve is the test ring]
  \label{thm:relStructureSectionsTop}
  \lean{AlgebraicGeometry.relStructureSectionsTop}
  \uses{def:relCurve, thm:isIso_appLE_snd}
  \leanok
  Let \(C\) be a smooth curve over \(k\) (proper, smooth of relative dimension one, geometrically
  irreducible), so that \(C\) is in particular geometrically reduced. For every commutative
  \(k\)-algebra \(R\), the degree-zero cohomology of the structure sheaf of the relative curve is
  the test ring itself:
  \[
    H^0\bigl(C_R, \struct{C_R/R}\bigr) \;=\; \Gamma(C_R, \top) \;\cong\; R,
  \]
  the isomorphism being pullback of global sections along the second projection
  \(\mathrm{pr}_2 : C_R \to \Spec R\). This is the degree-zero case of ``cohomology commutes with
  base change'', universal in the test ring \(R\).
\end{theorem}
\begin{proof}
  \leanok
  \uses{def:moduleKSheafHZero}
  Degree-zero cohomology is global sections, \(H^0(C_R, \struct{C_R/R}) \cong \Gamma(C_R, \top)\)
  (\ref{def:moduleKSheafHZero}, over the base ring \(R\)). Specialize the base-change isomorphism
  \ref{thm:isIso_appLE_snd} to the test object \(T = \Spec R\) and its affine open \(V = \top\):
  pullback along the second projection is an isomorphism
  \(\Gamma(\Spec R, \top) \xrightarrow{\sim} \Gamma(C_R, \mathrm{pr}_2^{-1}\top) = \Gamma(C_R, \top)\)
  (the hypotheses of \ref{thm:isIso_appLE_snd} --- \(C\) proper, geometrically irreducible and
  geometrically reduced --- hold for a smooth curve). Composing with the identification
  \(\Gamma(\Spec R, \top) \cong R\) gives \(\Gamma(C_R, \top) \cong R\).
\end{proof}

\section{The projective line: inputs for the finiteness argument}

This section collects the facts about \(\PP^1\) and about morphisms over \(\Spec k\) that
the finiteness argument of the next section consumes. It builds on the description of
\(\PP^1\), its standard charts \(U_0, U_1\), the chart coordinates
\(t_{ij} = X_j/X_i \in A_{(X_i)}\), and the section-ring identifications
\(\gamma_0, \gamma_1, \gamma_{01}\) developed in Chapter~\ref{chap:Curves}
(Definitions~\ref{def:P1_chartOpen}, \ref{def:P1_chartCoord},
\ref{def:P1_chartSectionsEquiv0}--\ref{def:P1_overlapSectionsEquiv}); the statements here
are recorded in this chapter because they are private inputs of the finiteness glue rather
than general theory of the projective line. Throughout, for a homogeneous \(f\) of positive
degree we write \(A_{(f)} \to \Gamma(\PP^1, D_+(f))\) for the canonical map from the
degree-zero localization to the sections over the basic open (an isomorphism).

\begin{lemma}[Morphisms over the base intertwine the structure algebra maps]
  \label{lem:appTop_map_appLE}
  \lean{AlgebraicGeometry.Scheme.Hom.appTop_map_appLE}
  \leanok
  Let \(k\) be a commutative ring, let \(X, Y\) be schemes over \(\Spec k\) and let
  \(\pi : X \to Y\) be a morphism compatible with the structure morphisms. For opens
  \(U \subseteq Y\) and \(V \subseteq X\) with \(V \subseteq \pi^{-1}(U)\), the two
  composite ring homomorphisms
  \[
    \Gamma(\Spec k, \top) \to \Gamma(Y, \top) \to \Gamma(Y, U)
      \xrightarrow{\ \pi^\sharp\ } \Gamma(X, V)
    \qquad\text{and}\qquad
    \Gamma(\Spec k, \top) \to \Gamma(X, \top) \to \Gamma(X, V)
  \]
  agree, where the unlabeled maps are induced by the structure morphisms and by
  restriction.
\end{lemma}
\begin{proof}
  \leanok
  Pulling back sections along \(\pi\) commutes with restriction, so the first composite is
  the map on sections induced by \(\pi\) composed with the structure morphism of \(Y\),
  restricted to \(V\); by compatibility this composite of morphisms is the structure
  morphism of \(X\).
\end{proof}

\begin{lemma}[Pulling back the structure algebra map]
  \label{lem:appLE_overAlgebraMap}
  \lean{AlgebraicGeometry.Scheme.Hom.appLE_overAlgebraMap}
  \uses{def:overAlgebraMap, lem:appTop_map_appLE}
  \leanok
  In the situation of \ref{lem:appTop_map_appLE}, for every \(r \in k\) the map
  \(\pi^\sharp : \Gamma(Y, U) \to \Gamma(X, V)\) sends the image of \(r\) under the
  structure algebra map of \(Y\) (Definition~\ref{def:overAlgebraMap}) to the image of
  \(r\) under the structure algebra map of \(X\).
\end{lemma}
\begin{proof}
  \leanok
  Evaluate the equality of ring homomorphisms of \ref{lem:appTop_map_appLE} at the section
  corresponding to \(r\).
\end{proof}

\begin{lemma}[The basic-open models lie over the base]
  \label{lem:P1_awayi_structureMap}
  \lean{AlgebraicGeometry.P1.awayι_structureMap}
  \uses{def:P1_structureMap}
  \leanok
  Let \(f \in k[X_0, X_1]\) be homogeneous of positive degree. The composite of the affine
  model \(\Spec A_{(f)} \to \PP^1\) of the basic open \(D_+(f)\) with the structure
  morphism \(\PP^1 \to \Spec k\) is the morphism induced by the algebra map
  \(k \to A_{(f)}\). (This generalizes \ref{lem:P1_chartIota_structureMap} from the two
  standard charts to all standard basic opens.)
\end{lemma}
\begin{proof}
  \leanok
  \uses{def:P1_gradeZero}
  As for \ref{lem:P1_chartIota_structureMap}: the canonical morphism
  \(\Proj A \to \Spec A_0\), restricted to the affine model of a basic open, is induced by
  the ring homomorphism \(A_0 \to A_{(f)}\), and composing with the isomorphism
  \(k \cong A_0\) of \ref{def:P1_gradeZero} gives the claim.
\end{proof}

\begin{lemma}[The basic-open inclusion factors through its affine model]
  \label{lem:P1_basicOpen_iota_eq}
  \lean{AlgebraicGeometry.P1.basicOpen_ι_eq}
  \uses{def:P1}
  \leanok
  Let \(f \in k[X_0, X_1]\) be homogeneous of positive degree. The inclusion
  \(D_+(f) \to \PP^1\) is the composite of the canonical identification
  \(D_+(f) \cong \Spec A_{(f)}\) with the affine model \(\Spec A_{(f)} \to \PP^1\).
\end{lemma}
\begin{proof}
  \leanok
  The affine model is by construction an open immersion with image \(D_+(f)\), and the
  canonical identification is the induced isomorphism onto the image; composing the two
  recovers the inclusion.
\end{proof}

\begin{lemma}[The structure algebra map on a basic open]
  \label{lem:P1_structureMap_appTop_awayToSection}
  \lean{AlgebraicGeometry.P1.structureMap_appTop_awayToSection}
  \uses{def:P1_structureMap, lem:P1_awayi_structureMap, lem:P1_basicOpen_iota_eq}
  \leanok
  Let \(f \in k[X_0, X_1]\) be homogeneous of positive degree. The composite
  \[
    \Gamma(\Spec k, \top) \to \Gamma(\PP^1, \top) \to \Gamma\bigl(\PP^1, D_+(f)\bigr)
  \]
  (structure morphism on global sections, then restriction) is the composite of the
  identification \(\Gamma(\Spec k, \top) \cong k\), the algebra map \(k \to A_{(f)}\), and
  the canonical map \(A_{(f)} \to \Gamma(\PP^1, D_+(f))\).
\end{lemma}
\begin{proof}
  \leanok
  By \ref{lem:P1_basicOpen_iota_eq} and \ref{lem:P1_awayi_structureMap}, the composite of
  the inclusion \(D_+(f) \to \PP^1\) with the structure morphism is the identification
  \(D_+(f) \cong \Spec A_{(f)}\) followed by the morphism induced by \(k \to A_{(f)}\).
  Taking global sections of this equality of scheme morphisms, and identifying the global
  sections of the open subscheme \(D_+(f)\) with \(\Gamma(\PP^1, D_+(f))\) --- under which
  the map on sections induced by the identification with \(\Spec A_{(f)}\) becomes the
  canonical map \(A_{(f)} \to \Gamma(\PP^1, D_+(f))\) --- yields the claim.
\end{proof}

\begin{lemma}[The distinguished fraction of a chart pair is the chart coordinate]
  \label{lem:P1_isLocalizationElem_X_eq}
  \lean{AlgebraicGeometry.P1.isLocalizationElem_X_eq}
  \uses{def:P1_chartCoord}
  \leanok
  For \(i, j \in \{0,1\}\), the canonical degree-zero fraction
  \(X_j^{\deg X_i} / X_i^{\deg X_j}\) attached to the pair \((X_i, X_j)\) of homogeneous
  elements equals the chart coordinate \(t_{ij} = X_j / X_i \in A_{(X_i)}\).
\end{lemma}
\begin{proof}
  \leanok
  Both degrees are one, and a first power is the element itself.
\end{proof}

\begin{theorem}[The chart coordinate cuts out the overlap]
  \label{thm:P1_basicOpen_chartCoord}
  \lean{AlgebraicGeometry.P1.basicOpen_awayToSection_chartCoord}
  \uses{def:P1_chartOpen, def:P1_chartCoord, lem:P1_basicOpen_iota_eq,
    lem:P1_isLocalizationElem_X_eq}
  \leanok
  For \(i, j \in \{0,1\}\), regard the chart coordinate \(t_{ij} \in A_{(X_i)}\) as a
  section of the structure sheaf over the chart \(U_i = D_+(X_i)\), via the canonical map
  \(A_{(X_i)} \to \Gamma(\PP^1, U_i)\). Then the basic open of this section in the scheme
  \(\PP^1\) --- the locus of \(U_i\) where it is invertible --- is the chart overlap:
  \[
    D(t_{ij}) \;=\; U_i \cap U_j .
  \]
\end{theorem}
\begin{proof}
  \leanok
  Both sides are open subsets of \(U_i\) (the basic open of a section over \(U_i\) is
  contained in \(U_i\)), so it suffices to show that they have the same preimage under the
  identification \(U_i \cong \Spec A_{(X_i)}\), and to transport the equality forward
  along this open immersion. No stalks are involved: everything is transported through
  the affine model.

  On the one hand, the preimage of the basic open \(D(t_{ij})\) is the basic open of
  \(\Spec A_{(X_i)}\) at the ring element \(t_{ij}\): under the identification of an
  affine open with the spectrum of its section ring, scheme-level basic opens of sections
  correspond to spectrum-level basic opens of ring elements.

  On the other hand, the preimage of \(U_j = D_+(X_j)\) under the affine model
  \(\Spec A_{(X_i)} \to \PP^1\) (through which the inclusion of \(U_i\) factors,
  \ref{lem:P1_basicOpen_iota_eq}) is the basic open at the canonical degree-zero fraction
  attached to the pair \((X_i, X_j)\): a homogeneous prime avoiding \(X_i\) avoids
  \(X_j\) exactly when the degree-zero fraction \(X_j/X_i\) is invertible in the
  corresponding localization. By \ref{lem:P1_isLocalizationElem_X_eq} this fraction is
  \(t_{ij}\), so the preimage of \(U_i \cap U_j\) is again the basic open at
  \(t_{ij}\).
\end{proof}

\begin{lemma}[The Laurent generators as restrictions from the charts]
  \label{lem:P1_laurent_generators}
  \lean{AlgebraicGeometry.P1.overlapSectionsEquiv_symm_T,
    AlgebraicGeometry.P1.overlapSectionsEquiv_symm_T_neg,
    AlgebraicGeometry.P1.overlapSectionsEquiv_symm_algebraMap}
  \uses{def:P1_overlapSectionsEquiv, def:P1_chartCoord}
  \leanok
  Under the identification
  \(\gamma_{01} : \Gamma(\PP^1, U_0 \cap U_1) \cong k[T, T^{-1}]\) of
  \ref{def:P1_overlapSectionsEquiv}:
  \begin{enumerate}
    \item \(T\) corresponds to the restriction to the overlap of the chart coordinate
      \(t_{01}\), regarded as a section over \(U_0\);
    \item \(T^{-1}\) corresponds to the restriction to the overlap of the chart
      coordinate \(t_{10}\), regarded as a section over \(U_1\);
    \item each scalar \(r \in k\) corresponds to the section over the overlap given by
      the image of \(r\) under the algebra map \(k \to A_{(X_0X_1)}\).
  \end{enumerate}
\end{lemma}
\begin{proof}
  \leanok
  \uses{lem:P1_res_left, lem:P1_res_right}
  (1) and (2): by \ref{lem:P1_res_left} (resp.\ \ref{lem:P1_res_right}), restriction from
  the chart \(U_0\) (resp.\ \(U_1\)) to the overlap corresponds, under the
  identifications \(\gamma_0, \gamma_1, \gamma_{01}\), to \(t \mapsto T\) (resp.\
  \(t \mapsto T^{-1}\)); apply this to the sections corresponding to the polynomial
  \(t\), which are the chart coordinates. (3): \(\gamma_{01}\) is built from the
  \(k\)-algebra isomorphism \(\theta\) of \ref{def:P1_overlapAlgEquiv}, hence matches the
  two algebra maps from \(k\).
\end{proof}

\section{Finiteness of the first cohomology}

Combining the two-affine-cover cokernel bridge with the two-lattice ladder of
Chapter~\ref{chap:Algebra} along a finite morphism to the projective line yields the
finiteness of \(H^1\). Classically this computation goes back to the explicit \v{C}ech
computation of the cohomology of projective space --- for \(\PP^1\) itself, the
description of \(H^1\) as the cokernel of the two-chart restriction-difference map into
the Laurent ring, spanned by the finitely many ``negative'' monomials --- extended to a
curve by pushing forward along a finite morphism; the two-lattice ladder is exactly the
abstraction of that monomial bookkeeping.

\begin{theorem}[Finiteness of \(H^1\) along a finite morphism to the projective line]
  \label{thm:moduleFinite_hModule_one_of_isFinite_toP1}
  \lean{AlgebraicGeometry.moduleFinite_hModule_one_of_isFinite_toP1}
  \uses{def:HModule, def:moduleKSheaf, def:P1_structureMap}
  \leanok
  \source{hartshorne-algebraic-geometry:page-0242, hartshorne-algebraic-geometry:page-0243}
  Let \(k\) be a field, \(X\) a scheme over \(\Spec k\), and \(\pi : X \to \PP^1\) a
  finite morphism over \(k\) (its composite with the structure morphism of \(\PP^1\) is
  the structure morphism of \(X\)). Then \(H^1\bigl(X, \struct{X/k}\bigr)\) is a finite
  \(k\)-module.
\end{theorem}
\begin{proof}
  \leanok
  \uses{lem:preimage_chart_affine, lem:preimage_chart_sup, lem:finite_app_overlap,
    lem:P1_chartOpen_inf, thm:TwoCover_h1CokEquiv, def:TwoCover_diff, def:TwoCover_H1Cok,
    def:P1_overlapSectionsEquiv, def:P1_chartCoord, lem:P1_laurent_generators,
    lem:appLE_overAlgebraMap, lem:P1_structureMap_appTop_awayToSection,
    thm:P1_basicOpen_chartCoord, cor:two_lattice_loc_pow}
  Write \(U_0, U_1\) for the standard charts of \(\PP^1\) and set
  \(V_i = \pi^{-1}(U_i)\).

  (a) \emph{The cover and the cokernel.} Since \(\pi\) is finite, hence affine, the
  \(V_i\) are affine opens of \(X\) (\ref{lem:preimage_chart_affine}), and they cover
  \(X\) (\ref{lem:preimage_chart_sup}). By the two-affine-cover bridge
  \ref{thm:TwoCover_h1CokEquiv},
  \[
    H^1\bigl(X, \struct{X/k}\bigr) \;\cong\;
      N \big/ \mathrm{range}(\mathrm{diff}),
    \qquad N := \Gamma(X, V_0 \cap V_1),
  \]
  with \(\mathrm{diff}(s_0, s_1) = s_0|_{V_0 \cap V_1} - s_1|_{V_0 \cap V_1}\)
  (\ref{def:TwoCover_diff}). It therefore suffices to prove that this quotient is a
  finite \(k\)-module.

  (b) \emph{The Laurent module structure.} Preimages commute with intersections, and
  \(U_0 \cap U_1 = D_+(X_0X_1)\) (\ref{lem:P1_chartOpen_inf}), so
  \(V_0 \cap V_1 = \pi^{-1}(U_0 \cap U_1)\). Make \(N\) an algebra over the Laurent
  polynomial ring via the composite
  \[
    k[T, T^{-1}] \xrightarrow{\ \gamma_{01}^{-1}\ } \Gamma(\PP^1, U_0 \cap U_1)
      \xrightarrow{\ \pi^\sharp\ } \Gamma\bigl(X, \pi^{-1}(U_0 \cap U_1)\bigr) = N
  \]
  (\ref{def:P1_overlapSectionsEquiv}). This structure is compatible with the \(k\)-module
  structure of \(N\): a scalar \(r \in k[T,T^{-1}]\) corresponds under \(\gamma_{01}\) to
  the image of \(r\) under the algebra map into the overlap ring
  (\ref{lem:P1_laurent_generators}(3)), which is the structure algebra map of \(\PP^1\)
  on \(U_0 \cap U_1\) (\ref{lem:P1_structureMap_appTop_awayToSection}), and \(\pi\) pulls
  the structure algebra map of \(\PP^1\) back to that of \(X\)
  (\ref{lem:appLE_overAlgebraMap}). Hence the \(k\)-module structure of \(N\) is the
  restriction of scalars of the \(k[T,T^{-1}]\)-structure along the canonical map
  \(k \to k[T,T^{-1}]\).

  (c) \emph{Module-finiteness over the Laurent ring.} Since \(\pi\) is finite and
  \(U_0 \cap U_1\) is an affine open of \(\PP^1\), the map
  \(\pi^\sharp : \Gamma(\PP^1, U_0 \cap U_1) \to N\) is module-finite
  (\ref{lem:finite_app_overlap}); composing with the isomorphism \(\gamma_{01}^{-1}\),
  \(N\) is a finite \(k[T,T^{-1}]\)-module.

  (d) \emph{The two lattices, and the range of the difference map.} Let
  \(Q_0, Q_1 \subseteq N\) be the images of the (\(k\)-linear) restriction maps
  \(\Gamma(X, V_0) \to N\) and \(\Gamma(X, V_1) \to N\). Then
  \(\mathrm{range}(\mathrm{diff}) = Q_0 + Q_1\): every value
  \(s_0| - s_1|\) of \(\mathrm{diff}\) is the sum of \(s_0| \in Q_0\) and
  \(-(s_1|) = (-s_1)| \in Q_1\); conversely \(s_0| = \mathrm{diff}(s_0, 0) \) and
  \(s_1| = \mathrm{diff}(0, -s_1)\) exhibit every element of \(Q_0\) and of \(Q_1\) as a
  value of \(\mathrm{diff}\) --- the sign is absorbed because \(Q_1\) is the image of a
  linear map.

  (e) \emph{Lattice stability.} Let
  \(\tau_0 = \pi^\sharp(t_{01}) \in \Gamma(X, V_0)\) and
  \(\tau_1 = \pi^\sharp(t_{10}) \in \Gamma(X, V_1)\) be the pullbacks of the chart
  coordinates, regarded as sections over the charts
  (Definition~\ref{def:P1_chartCoord}). By \ref{lem:P1_laurent_generators}(1) and the
  compatibility of \(\pi^\sharp\) with restriction, \(T\) acts on \(N\) as
  multiplication by \(\tau_0|_{V_0 \cap V_1}\); symmetrically \(T^{-1}\) acts as
  multiplication by \(\tau_1|_{V_0 \cap V_1}\) (\ref{lem:P1_laurent_generators}(2)).
  Hence \(Q_0\) is stable under multiplication by \(T\): for \(s \in \Gamma(X, V_0)\),
  \[
    T \cdot s|_{V_0 \cap V_1} \;=\; \tau_0|\cdot s| \;=\; (\tau_0 s)|_{V_0 \cap V_1}
      \;\in\; Q_0,
  \]
  and symmetrically \(Q_1\) is stable under multiplication by \(T^{-1}\).

  (f) \emph{Absorption up to powers.} The overlap is a basic open inside each chart
  preimage: pulling back the identity \(D(t_{ij}) = U_i \cap U_j\) of
  \ref{thm:P1_basicOpen_chartCoord} along \(\pi\) (basic opens of sections pull back to
  basic opens of the pulled-back sections) gives
  \[
    V_0 \cap V_1 = D(\tau_0) \subseteq V_0,
    \qquad
    V_0 \cap V_1 = D(\tau_1) \subseteq V_1 .
  \]
  Since \(V_0\) is affine and \(V_0 \cap V_1\) is its basic open defined by \(\tau_0\),
  restriction identifies \(N\) with the localization of \(\Gamma(X, V_0)\) at the powers
  of \(\tau_0\) (quasi-coherence of the structure sheaf). Hence for every \(n \in N\)
  there are \(m \ge 0\) and \(a \in \Gamma(X, V_0)\) with
  \(n \cdot (\tau_0|)^m = a|\), i.e.\ \(T^m \cdot n \in Q_0\). Symmetrically, every
  \(n \in N\) satisfies \(T^{-m} \cdot n \in Q_1\) for some \(m \ge 0\).

  (g) \emph{The ladder.} By (b), (c), (e) and (f), the module \(N\) over
  \(k[T,T^{-1}]\) with the \(k\)-submodules \(Q_0, Q_1\) satisfies all hypotheses of the
  two-lattice finiteness theorem in its localization form via powers of the generator
  (\ref{cor:two_lattice_loc_pow}); hence \(N/(Q_0 + Q_1)\) is a finite \(k\)-module. By
  (d) this quotient is \(N/\mathrm{range}(\mathrm{diff})\), and by (a) it is
  \(H^1\bigl(X, \struct{X/k}\bigr)\).
\end{proof}

\begin{theorem}[Finiteness of \(H^1\) of the curve]
  \label{thm:moduleFinite_hModule_one}
  \lean{AlgebraicGeometry.moduleFinite_hModule_one}
  \uses{def:HModule, def:moduleKSheaf}
  \leanok
  Let \(k\) be a field and \(C\) a smooth curve over \(k\) (proper, smooth of relative
  dimension one, geometrically irreducible). Then
  \(H^1\bigl(C, \struct{C/k}\bigr)\) is a finite \(k\)-module; equivalently, it is a
  finite-dimensional \(k\)-vector space.
\end{theorem}
\begin{proof}
  \leanok
  \uses{thm:exists_finite_toP1, thm:moduleFinite_hModule_one_of_isFinite_toP1}
  By \ref{thm:exists_finite_toP1} there is a finite morphism \(\pi : C \to \PP^1\) over
  \(k\); apply \ref{thm:moduleFinite_hModule_one_of_isFinite_toP1}.
\end{proof}

\begin{remark}[The genus is the dimension of a finite-dimensional space]
  \label{rem:genus_finite}
  \uses{def:genus, thm:moduleFinite_hModule_one}
  \leanok
  For a smooth curve \(C\) over \(k\), the genus of Definition~\ref{def:genus} is the
  \(k\)-dimension of \(H^1(C, \struct C)\), which by
  \ref{thm:moduleFinite_hModule_one} is a finite-dimensional \(k\)-vector space; the
  genus is therefore a well-defined natural number, attained as the cardinality of any
  basis.
\end{remark}

\section{The cocycle-glued twisted sheaf on a two-cover}
\label{sec:twisted_sheaf}

Let \(k\) be a commutative ring and \(X\) a scheme over \(\Spec k\), its structure sheaf
viewed as a sheaf of \(k\)-modules (\ref{def:moduleKSheaf}). Fix two opens
\(V_0, V_1 \subseteq X\) and a \emph{unit cocycle}: an invertible section
\(g \in \Gamma(X, V_0 \cap V_1)^{\times}\). (On a two-chart cover a \v{C}ech
\(1\)-cocycle has a single component and the cocycle condition on triple overlaps is
vacuous.) This section glues from \(g\) a sheaf \(F_g\) of \(k\)-modules --- for
\(V_0 \cup V_1 = X\), the line bundle with transition function \(g\), trivialised on each
chart --- and computes its degree-one cohomology as a two-cover cokernel, by transporting
the quasi-coherence packaging of Section~\ref{sec:qcoh_vanishing} through the chart
trivialisations.

\begin{definition}[Transport of the packaging along a sections-level trivialisation]
  \label{def:qcohOn_transport}
  \lean{AlgebraicGeometry.Scheme.QcohOn.ofSectionsEquiv}
  \uses{def:qcohOn}
  \leanok
  Let \(F\) and \(G\) be sheaves of \(k\)-modules on the small Zariski site of \(X\), let
  \(U\) be an open, and suppose given \(k\)-linear equivalences
  \(e_W : F(W) \cong G(W)\) for every open \(W \le U\), commuting with section
  restriction. If \(G\) carries a quasi-coherence packaging on \(U\)
  (\ref{def:qcohOn}), then so does \(F\), with the conjugated action
  \(r \cdot m := e_W^{-1}\bigl(r \cdot e_W(m)\bigr)\).
\end{definition}
\begin{proof}
  \leanok
  Each (P1) law for \(F\) follows by conjugating the corresponding law for \(G\) through
  the equivalences, and naturality under restriction combines the (P1) naturality for
  \(G\) with the compatibility of \(e\) with restriction. For (P2): given
  \(c \in F(V \cap D(g))\) on an affine \(V \le U\), denominator clearing for \(G\)
  applied to \(e(c)\) produces \(N\) and \(e' \in G(V)\) with
  \(e'|_{V \cap D(g)} = g^N \cdot e(c)\); then \(e_V^{-1}(e')\) clears the denominator
  for \(c\), since \(e\) commutes with restriction and the action is conjugated. Defect
  annihilation transports the same way.
\end{proof}

\begin{definition}[Twisted sections of a unit cocycle]
  \label{def:twistSections}
  \lean{AlgebraicGeometry.twistDefect, AlgebraicGeometry.twistSubmodule}
  \uses{def:moduleKSheaf}
  \leanok
  For an open \(W \subseteq X\), the \emph{defect map} of the cocycle \(g\) over \(W\) is
  the \(k\)-linear map
  \[
    \Gamma(X, W \cap V_0) \times \Gamma(X, W \cap V_1) \longrightarrow
      \Gamma(X, W \cap V_0 \cap V_1), \qquad
    (s_0, s_1) \longmapsto s_0| - g \cdot s_1|,
  \]
  where both restrictions are to \(W \cap V_0 \cap V_1\) and \(g\) acts through its
  restriction there. The module of \emph{twisted sections} \(F_g(W)\) is the kernel of
  the defect map: the pairs \((s_0, s_1)\) that match through \(g\) on the overlap.
\end{definition}

\begin{lemma}[Twisted sections form a presheaf]
  \label{lem:twistSections_presheaf}
  \lean{AlgebraicGeometry.twistRes, AlgebraicGeometry.twistSubmodule_res,
    AlgebraicGeometry.twistPresheaf}
  \uses{def:twistSections}
  \leanok
  For \(W' \le W\), componentwise restriction carries \(F_g(W)\) into \(F_g(W')\), and
  the resulting \(k\)-linear maps are functorial in the inclusion; so \(W \mapsto F_g(W)\)
  is a presheaf of \(k\)-modules.
\end{lemma}
\begin{proof}
  \leanok
  Restrict the matching identity \(s_0| = g \cdot s_1|\) from \(W \cap V_0 \cap V_1\) to
  \(W' \cap V_0 \cap V_1\): restriction of the structure sheaf is a ring homomorphism
  commuting with further restriction, so the restricted pair again matches.
  Functoriality is functoriality of restriction in each component.
\end{proof}

\begin{theorem}[The twisted sheaf]
  \label{thm:twistSheaf}
  \lean{AlgebraicGeometry.isSheaf_twistPresheaf, AlgebraicGeometry.twistSheaf}
  \uses{lem:twistSections_presheaf}
  \leanok
  The presheaf \(F_g\) of twisted sections is a sheaf of \(k\)-modules on the small
  Zariski site of \(X\).
\end{theorem}
\begin{proof}
  \leanok
  Let \(\{U_i\}\) be an open cover of \(W\) and \((s_0^i, s_1^i) \in F_g(U_i)\) sections
  agreeing on overlaps. The first components are sections of the structure sheaf over the
  cover \(\{U_i \cap V_0\}\) of \(W \cap V_0\), compatible on its overlaps; by the sheaf
  property of the structure sheaf they glue to a unique \(t_0 \in \Gamma(X, W \cap V_0)\),
  and likewise the second components glue to a unique \(t_1 \in \Gamma(X, W \cap V_1)\).
  The matching \(t_0| = g \cdot t_1|\) on \(W \cap V_0 \cap V_1\) may be checked on the
  cover \(\{U_i \cap V_0 \cap V_1\}\), where it restricts to the matching of
  \((s_0^i, s_1^i)\); by separation of the structure sheaf it holds on the union. So
  \((t_0, t_1) \in F_g(W)\) restricts to every \((s_0^i, s_1^i)\), and it is the unique
  such element because each of its components is determined by separation.
\end{proof}

\begin{definition}[Chart trivialisations of the twisted sheaf]
  \label{def:twistTriv}
  \lean{AlgebraicGeometry.twistTriv₀, AlgebraicGeometry.twistTriv₁}
  \uses{def:twistSections}
  \leanok
  Let \(W \le V_0\) (so \(W \cap V_0 = W\) and \(W \cap V_0 \cap V_1 = W \cap V_1\)).
  Projection to the first component is a \(k\)-linear equivalence
  \[
    \mathrm{triv}_0 : F_g(W) \;\cong\; \Gamma(X, W), \qquad (s_0, s_1) \mapsto s_0,
  \]
  with inverse \(s \mapsto (s,\ g^{-1} \cdot s|_{W \cap V_1})\). Symmetrically, for
  \(W \le V_1\), projection to the second component is a \(k\)-linear equivalence
  \(\mathrm{triv}_1 : F_g(W) \cong \Gamma(X, W)\) with inverse
  \(s \mapsto (g \cdot s|_{W \cap V_0},\ s)\).
\end{definition}
\begin{proof}
  \leanok
  For \(W \le V_0\) the matching relation of a pair \((s_0, s_1) \in F_g(W)\) reads
  \(s_0| = g \cdot s_1\) on \(W \cap V_1\), i.e.\ \(s_1 = g^{-1} \cdot s_0|\): the second
  component is rebuilt from the first, so projection and the displayed inverse are
  mutually inverse; the pair \((s, g^{-1} \cdot s|)\) matches by construction. Both maps
  are \(k\)-linear, multiplication by the fixed sections \(g^{\pm 1}\) being
  \(k\)-linear. The second chart is symmetric.
\end{proof}

\begin{lemma}[The trivialisations commute with restriction]
  \label{lem:twistTriv_res}
  \lean{AlgebraicGeometry.twistTriv₀_res, AlgebraicGeometry.twistTriv₁_res}
  \uses{def:twistTriv, lem:twistSections_presheaf}
  \leanok
  For \(W' \le W \le V_0\), trivialising after restricting equals restricting after
  trivialising: \(\mathrm{triv}_0(p|_{W'}) = (\mathrm{triv}_0\, p)|_{W'}\); likewise on
  the second chart.
\end{lemma}
\begin{proof}
  \leanok
  Both sides are the first (resp.\ second) component of \(p\) restricted to \(W'\);
  componentwise restriction is functorial.
\end{proof}

\begin{theorem}[The twisted sheaf is quasi-coherently packaged on each chart]
  \label{thm:twistSheaf_qcohOn}
  \lean{AlgebraicGeometry.twistSheaf.instQcohOn₀, AlgebraicGeometry.twistSheaf.instQcohOn₁}
  \uses{thm:twistSheaf, def:qcohOn_transport, def:twistTriv, lem:twistTriv_res,
    thm:moduleKSheaf_qcohOn}
  \leanok
  The twisted sheaf \(F_g\) carries a quasi-coherence packaging on \(V_0\) and on
  \(V_1\).
\end{theorem}
\begin{proof}
  \leanok
  The chart trivialisations give \(k\)-linear equivalences
  \(F_g(W) \cong \Gamma(X, W)\) for every \(W \le V_0\) (resp.\ \(W \le V_1\)),
  compatible with restriction (\ref{lem:twistTriv_res}); the structure sheaf is packaged
  on every open (\ref{thm:moduleKSheaf_qcohOn}); transport the packaging along the
  trivialisation (\ref{def:qcohOn_transport}).
\end{proof}

\begin{theorem}[Twisted affine vanishing]
  \label{thm:twistSheaf_affine_vanishing}
  \lean{AlgebraicGeometry.subsingleton_hModule'_twistSheaf_one₀,
    AlgebraicGeometry.subsingleton_hModule'_twistSheaf_one₁}
  \uses{def:HModule', thm:twistSheaf_qcohOn, thm:hModule1_vanishing_qcoh}
  \leanok
  If \(V_0\) (resp.\ \(V_1\)) is an affine open of \(X\), then
  \(H'^1(V_0, F_g) = 0\) (resp.\ \(H'^1(V_1, F_g) = 0\)).
\end{theorem}
\begin{proof}
  \leanok
  The twisted affine vanishing theorem \ref{thm:hModule1_vanishing_qcoh}, applied to the
  packaging of \ref{thm:twistSheaf_qcohOn}.
\end{proof}

\begin{definition}[The two-cover \(H^1\) carrier of the twisted sheaf]
  \label{def:twistTwoCoverH1}
  \lean{AlgebraicGeometry.twistTwoCoverH1}
  \uses{thm:twistSheaf, def:twoCoverSquare, thm:twistSheaf_affine_vanishing,
    thm:twoCoverH1LinearEquiv}
  \leanok
  Let \(V_0, V_1\) be affine opens with \(V_0 \cup V_1 = X\). Then the degree-one
  cohomology of the twisted sheaf is the two-cover cokernel of twisted sections: a
  \(k\)-linear isomorphism
  \[
    H^1(X, F_g) \;\cong\;
      F_g(V_0 \cap V_1) \big/ \mathrm{range}\bigl((p_0, p_1) \mapsto
        p_0|_{V_0 \cap V_1} - p_1|_{V_0 \cap V_1}\bigr),
  \]
  the restriction-difference map being that of the two-cover square
  (\ref{def:twoCoverSquare}) with coefficients in \(F_g\). (The overlap sections
  \(F_g(V_0 \cap V_1)\) are further identified with \(\Gamma(X, V_0 \cap V_1)\) by the
  chart trivialisation at \(V_0 \cap V_1 \le V_0\), \ref{def:twistTriv}.)
\end{definition}
\begin{proof}
  \leanok
  The two charts are affine, so \(H'^1(V_i, F_g) = 0\)
  (\ref{thm:twistSheaf_affine_vanishing}); apply the two-cover \(H^1\) computation with
  general coefficients (\ref{thm:twoCoverH1LinearEquiv}) to \(F = F_g\).
\end{proof}

\begin{definition}[The relative twisted sheaf]
  \label{def:relTwistSheaf}
  \lean{AlgebraicGeometry.relUnitCocycle, AlgebraicGeometry.relTwistSheaf}
  \uses{def:relCurve, def:relCover, lem:relCover_inf, def:twistSections, thm:twistSheaf}
  \leanok
  Let \(k\) be a field, \(C \in \Over(\Spec k)\), \(R\) a commutative \(k\)-algebra, and
  \(D = (V_0, V_1)\) an affine two-chart cover of \(C\). A unit cocycle
  \(g_k \in \Gamma(C, V_0 \cap V_1)^{\times}\) pulls back along the first projection to a
  unit cocycle on the base-changed cover of the relative curve:
  \(\mathrm{pr}_1^{\sharp}(g_k) \in
  \Gamma\bigl(C_R,\ \mathrm{pr}_1^{-1}V_0 \cap \mathrm{pr}_1^{-1}V_1\bigr)^{\times}\)
  (pullback of sections is a ring homomorphism, so it preserves units; the overlap of the
  base-changed charts is the preimage of \(V_0 \cap V_1\) by \ref{lem:relCover_inf}). The
  \emph{relative twisted sheaf} of a unit cocycle \(g\) on the base-changed cover is the
  twisted sheaf \(F_g\) on \(C_R\), a sheaf of \(R\)-modules.
\end{definition}

\begin{definition}[The relative twisted \(H^1\) carrier]
  \label{def:relTwistTwoCoverH1}
  \lean{AlgebraicGeometry.relTwistTwoCoverH1}
  \uses{def:relTwistSheaf, def:twistTwoCoverH1, lem:relCover_charts}
  \leanok
  In the setting of \ref{def:relTwistSheaf}, the degree-one cohomology of the relative
  twisted sheaf is the explicit two-cover cokernel of \(R\)-modules: an \(R\)-linear
  isomorphism
  \[
    H^1(C_R, F_g) \;\cong\;
      F_g\bigl(\mathrm{pr}_1^{-1}V_0 \cap \mathrm{pr}_1^{-1}V_1\bigr) \big/
        \mathrm{range}(\mathrm{diff}) .
  \]
\end{definition}
\begin{proof}
  \leanok
  The base-changed charts are affine and cover \(C_R\) (\ref{lem:relCover_charts});
  apply \ref{def:twistTwoCoverH1} over the base ring \(R\).
\end{proof}

\section{Base change of the two-term complex and of \(H^1\) on the relative curve}
\label{sec:relative_h1_base_change}

Fix a field \(k\), an object \(C \in \Over(\Spec k)\), and a homomorphism \(R \to R'\) of
commutative \(k\)-algebras --- a change of affine test rings. This section identifies the
sections of the relative curve as the free base change of the sections of \(C\), builds
the comparison morphism \(C_{R'} \to C_R\), base-changes the two-term two-cover complex
of Section~\ref{sec:mv-in-text} along \(R \to R'\) termwise, and concludes that the
degree-one cohomology of the structure sheaf of the relative curve commutes with
\emph{every} change of test ring --- no flatness of \(R \to R'\) is needed, because the
two-cover \(H^1\) is a cokernel and cokernels commute with arbitrary base change
(Chapter~\ref{chap:Algebra}).

\begin{theorem}[Sections of the relative curve are the base change of sections]
  \label{thm:sectionsBaseChange}
  \lean{AlgebraicGeometry.Over.sectionsBaseChange, AlgebraicGeometry.Over.sectionsBaseChange_tmul}
  \uses{def:relCurve}
  \leanok
  Let \(k\) be a field, \(C \in \Over(\Spec k)\), \(R\) a commutative \(k\)-algebra, and
  \(V \subseteq C\) a quasi-compact quasi-separated open (for instance any affine open,
  or a finite intersection or union of affine opens of a quasi-separated scheme). Write
  \(V_R := \mathrm{pr}_1^{-1} V \subseteq C_R\). Then there is a ring isomorphism
  \[
    \Gamma(C, V) \otimes_k R \;\cong\; \Gamma(C_R, V_R), \qquad
    s \otimes a \longmapsto \mathrm{pr}_1^{\sharp}(s) \cdot \mathrm{pr}_2^{\sharp}(a),
  \]
  where \(\mathrm{pr}_1^{\sharp}\) is pullback of sections along the first projection,
  and \(\mathrm{pr}_2^{\sharp}(a)\) is the pullback along the second projection of the
  global section \(a \in R = \Gamma(\Spec R, \top)\). (The \(k\)-algebra structure on
  \(\Gamma(C, V)\) is restriction of scalars along the structure morphism.)
\end{theorem}
\begin{proof}
  \leanok
  The displayed assignment is a well-defined \(k\)-algebra map \(\chi_V\) out of the
  tensor product: both pullbacks are \(k\)-algebra homomorphisms into the commutative
  ring \(\Gamma(C_R, V_R)\), and the tensor product of commutative \(k\)-algebras is
  their pushout.

  \emph{Step 1: \(V\) affine.} The open \(V_R = \mathrm{pr}_1^{-1}(V)\) of the fibre
  product \(C_R = C \times_{\Spec k} \Spec R\) is the fibre product
  \(V \times_{\Spec k} \Spec R\) of the open subscheme \(V\), a fibre product of affine
  schemes, hence the affine scheme \(\Spec(\Gamma(V) \otimes_k R)\): \(\Spec\) converts
  the pushout of commutative rings into the fibre product of affine schemes. Under this
  identification \(\chi_V\) is the canonical isomorphism.

  \emph{Step 2: \(V\) quasi-compact quasi-separated.} Choose a finite affine cover
  \(V = U_1 \cup \dots \cup U_n\) (quasi-compactness); each intersection
  \(U_i \cap U_j\) is quasi-compact (quasi-separatedness), so it in turn has a finite
  affine cover \(\{W^{ij}_l\}\). The sheaf axiom writes \(\Gamma(C, V)\) as the
  equaliser of the two restriction maps between the finite products
  \(\prod_i \Gamma(U_i)\) and \(\prod_{i,j,l} \Gamma(W^{ij}_l)\). The preimages
  \(\mathrm{pr}_1^{-1} U_i\) form an affine cover of \(V_R\) (Step 1) whose pairwise
  intersections are covered by the affine \(\mathrm{pr}_1^{-1} W^{ij}_l\) (preimage
  commutes with unions and intersections), so \(\Gamma(C_R, V_R)\) is the equaliser of
  the corresponding restriction maps for the preimage cover. Since \(R\) is a free,
  hence flat, \(k\)-module (\(k\) is a field), the functor \(- \otimes_k R\) is exact
  and commutes with finite products, so it carries the first equaliser diagram to an
  equaliser diagram. The maps \(\chi\) on the terms commute with the two restriction
  maps (pullback commutes with restriction) and are isomorphisms on the product terms by
  Step 1; hence the induced map on equalisers, which is \(\chi_V\), is an isomorphism.
\end{proof}

\begin{definition}[Relative sections as a free \(R\)-module base change]
  \label{def:relSectionsBaseChange}
  \lean{AlgebraicGeometry.relPullbackSection, AlgebraicGeometry.relSectionsBaseChange,
    AlgebraicGeometry.relSectionsBaseChange_tmul,
    AlgebraicGeometry.Over.sectionsBaseChange_one_tmul_overAlgebraMap}
  \uses{thm:sectionsBaseChange}
  \leanok
  In the setting of \ref{thm:sectionsBaseChange}, endow \(\Gamma(C_R, V_R)\) with the
  \(R\)-module structure \(r \cdot y := \mathrm{pr}_2^{\sharp}(r) \cdot y\) given by
  pullback along the structure morphism \(\mathrm{pr}_2 : C_R \to \Spec R\). The
  isomorphism of \ref{thm:sectionsBaseChange}, read with the tensor factors flipped, is
  an \(R\)-linear equivalence
  \[
    R \otimes_k \Gamma(C, V) \;\cong\; \Gamma(C_R, V_R), \qquad
    a \otimes s \longmapsto \mathrm{pr}_2^{\sharp}(a) \cdot \mathrm{pr}_1^{\sharp}(s),
  \]
  with \(R\) acting on the left tensor factor; in particular
  \(1 \otimes r \mapsto \mathrm{pr}_2^{\sharp}(r)\).
\end{definition}
\begin{proof}
  \leanok
  Only \(R\)-linearity needs comment: on a pure tensor,
  \(r \cdot (a \otimes s) = (ra) \otimes s\) maps to
  \(\mathrm{pr}_2^{\sharp}(ra) \cdot \mathrm{pr}_1^{\sharp}(s) =
  \mathrm{pr}_2^{\sharp}(r) \cdot \bigl(\mathrm{pr}_2^{\sharp}(a) \cdot
  \mathrm{pr}_1^{\sharp}(s)\bigr)\), because \(\mathrm{pr}_2^{\sharp}\) is a ring
  homomorphism --- and the right-hand side is \(r\) acting on the image.
\end{proof}

\begin{definition}[The comparison morphism of relative curves]
  \label{def:relCurveMap}
  \lean{AlgebraicGeometry.overSpecMap, AlgebraicGeometry.relCurveMap}
  \uses{def:relCurve}
  \leanok
  Let \(\varphi : R \to R'\) be a homomorphism of commutative \(k\)-algebras. Applying
  \(\Spec\) gives a morphism \(\Spec R' \to \Spec R\) over \(\Spec k\), and the induced
  morphism of fibre products
  \[
    \rho \;=\; \mathrm{id}_C \times \Spec(\varphi) \;:\; C_{R'} \longrightarrow C_R
  \]
  is the \emph{comparison morphism} between the relative curves.
\end{definition}

\begin{lemma}[Projection identities of the comparison morphism]
  \label{lem:relCurveMap_proj}
  \lean{AlgebraicGeometry.relCurveMap_fst, AlgebraicGeometry.relCurveMap_snd,
    AlgebraicGeometry.relCurveMap_preimage}
  \uses{def:relCurveMap}
  \leanok
  The comparison morphism satisfies
  \(\mathrm{pr}_1 \circ \rho = \mathrm{pr}_1\) and
  \(\mathrm{pr}_2 \circ \rho = \Spec(\varphi) \circ \mathrm{pr}_2\); consequently
  \(\rho^{-1}(V_R) = V_{R'}\) for every open \(V \subseteq C\).
\end{lemma}
\begin{proof}
  \leanok
  The first two identities are the compatibility of the induced morphism of fibre
  products with the two projections; the third follows from the first by taking
  preimages of \(V\).
\end{proof}

\begin{definition}[The sections comparison map]
  \label{def:relSectionsMap}
  \lean{AlgebraicGeometry.relSectionsMap}
  \uses{def:relCurveMap, lem:relCurveMap_proj}
  \leanok
  For \(V \subseteq C\) open, pullback of sections along \(\rho\) (allowed since
  \(\rho^{-1}(V_R) = V_{R'}\), \ref{lem:relCurveMap_proj}) is a ring homomorphism
  \[
    \rho^{\sharp} \;:\; \Gamma(C_R, V_R) \longrightarrow \Gamma(C_{R'}, V_{R'}) .
  \]
\end{definition}

\begin{lemma}[Governing identities of the sections comparison map]
  \label{lem:relSectionsMap_props}
  \lean{AlgebraicGeometry.relSectionsMap_pullback,
    AlgebraicGeometry.relSectionsMap_overAlgebraMap,
    AlgebraicGeometry.relSectionsMap_smul, AlgebraicGeometry.relSectionsMap_resHom}
  \uses{def:relSectionsMap, lem:relCurveMap_proj, def:relSectionsBaseChange}
  \leanok
  The sections comparison map \(\rho^{\sharp}\):
  \begin{enumerate}
    \item commutes with the curve pullbacks:
      \(\rho^{\sharp}(\mathrm{pr}_1^{\sharp} s) = \mathrm{pr}_1^{\sharp} s\) for
      \(s \in \Gamma(C, V)\);
    \item intertwines the structure pullbacks:
      \(\rho^{\sharp}(\mathrm{pr}_2^{\sharp} r) = \mathrm{pr}_2^{\sharp}(\varphi(r))\)
      for \(r \in R\);
    \item hence is semilinear along \(\varphi\):
      \(\rho^{\sharp}(r \cdot y) = \varphi(r) \cdot \rho^{\sharp}(y)\) for the module
      structures of \ref{def:relSectionsBaseChange};
    \item commutes with restriction along \(W \le V\).
  \end{enumerate}
\end{lemma}
\begin{proof}
  \leanok
  (1) is pullback along \(\mathrm{pr}_1 \circ \rho = \mathrm{pr}_1\)
  (\ref{lem:relCurveMap_proj}); (2) is pullback along
  \(\mathrm{pr}_2 \circ \rho = \Spec(\varphi) \circ \mathrm{pr}_2\), and the pullback of
  the global section \(r\) along \(\Spec(\varphi)\) is \(\varphi(r)\); (3) is (2)
  together with multiplicativity of \(\rho^{\sharp}\); (4) is functoriality of pullback
  of sections with respect to inclusions of opens.
\end{proof}

\begin{theorem}[CBC-1: termwise base change of relative sections]
  \label{thm:relTermBaseChange}
  \lean{AlgebraicGeometry.relTermBaseChange, AlgebraicGeometry.relTermBaseChange_tmul}
  \uses{def:relSectionsBaseChange, lem:relSectionsMap_props}
  \leanok
  Let \(R \to R'\) be a homomorphism of test rings and \(V \subseteq C\) a quasi-compact
  quasi-separated open. There is an \(R'\)-linear equivalence
  \[
    R' \otimes_R \Gamma(C_R, V_R) \;\cong\; \Gamma(C_{R'}, V_{R'}), \qquad
    r' \otimes y \longmapsto r' \cdot \rho^{\sharp}(y) .
  \]
\end{theorem}
\begin{proof}
  \leanok
  Compose the free identifications of \ref{def:relSectionsBaseChange} over \(R\) and
  over \(R'\) with the tensor cancellation
  \(R' \otimes_R (R \otimes_k M) \cong R' \otimes_k M\):
  \[
    R' \otimes_R \Gamma(C_R, V_R) \;\cong\; R' \otimes_R \bigl(R \otimes_k \Gamma(C, V)\bigr)
      \;\cong\; R' \otimes_k \Gamma(C, V) \;\cong\; \Gamma(C_{R'}, V_{R'}) .
  \]
  For the formula, both sides are additive in \(y\) and \(R'\)-linear in \(r'\), so it
  suffices to treat \(y = \mathrm{pr}_2^{\sharp}(a) \cdot \mathrm{pr}_1^{\sharp}(s)\)
  (the images of pure tensors, which span). Chasing \(r' \otimes y\) through the
  composite gives \(\mathrm{pr}_2^{\sharp}\bigl(r' \varphi(a)\bigr) \cdot
  \mathrm{pr}_1^{\sharp}(s) = r' \cdot \bigl(\mathrm{pr}_2^{\sharp}(\varphi(a)) \cdot
  \mathrm{pr}_1^{\sharp}(s)\bigr)\), which by (1) and (2) of
  \ref{lem:relSectionsMap_props} is \(r' \cdot \rho^{\sharp}(y)\).
\end{proof}

\begin{lemma}[CBC-1: the two-cover complex base-changes]
  \label{lem:relDiffBaseChange}
  \lean{AlgebraicGeometry.relDiffDomBaseChange, AlgebraicGeometry.relDiffCodBaseChange,
    AlgebraicGeometry.relDiffBaseChange, AlgebraicGeometry.relDiffBaseChange_range}
  \uses{thm:relTermBaseChange, def:TwoCover_diff, def:relCover, lem:relCover_charts,
    lem:relCover_inf, lem:relSectionsMap_props}
  \leanok
  Let \(D = (V_0, V_1)\) be an affine two-chart cover of \(C\), and write
  \(\mathrm{diff}_R\) for the two-cover restriction-difference map
  (\ref{def:TwoCover_diff}) of \(C_R\) with its base-changed cover
  (\ref{def:relCover}). Under the term equivalences of \ref{thm:relTermBaseChange} ---
  the product of the two chart equivalences on the domain, the overlap equivalence on the
  codomain --- the base change \(\mathrm{diff}_R \otimes_R R'\) corresponds to
  \(\mathrm{diff}_{R'}\): the evident square commutes. In particular the overlap term
  equivalence carries \(\mathrm{range}(\mathrm{diff}_R \otimes_R R')\) onto
  \(\mathrm{range}(\mathrm{diff}_{R'})\).
\end{lemma}
\begin{proof}
  \leanok
  Both composites are additive, so it suffices to compare them on
  \(r' \otimes (s_0, s_1)\). Along the top and right, the result is
  \(r' \cdot \rho^{\sharp}\bigl(s_0|_{\cap} - s_1|_{\cap}\bigr)\); along the left and
  bottom, it is \(\mathrm{diff}_{R'}\bigl(r' \cdot \rho^{\sharp} s_0,\ r' \cdot
  \rho^{\sharp} s_1\bigr) = \bigl(r' \cdot \rho^{\sharp} s_0\bigr)|_{\cap} -
  \bigl(r' \cdot \rho^{\sharp} s_1\bigr)|_{\cap}\). These agree because
  \(\rho^{\sharp}\) commutes with restriction (\ref{lem:relSectionsMap_props}(4)) and
  the \(R'\)-action commutes with restriction (the structure pullback restricts to the
  structure pullback). The range statement follows because the domain equivalence is
  surjective.
\end{proof}

\begin{theorem}[CBC-2, cokernel form: the two-cover \(H^1\) carrier base-changes]
  \label{thm:relH1CokBaseChange}
  \lean{AlgebraicGeometry.relH1CokBaseChange, AlgebraicGeometry.relH1CokBaseChange_tmul_mk}
  \uses{thm:quot_range_baseChange, lem:relDiffBaseChange, def:TwoCover_H1Cok}
  \leanok
  In the setting of \ref{lem:relDiffBaseChange} there is an \(R'\)-linear equivalence
  \[
    R' \otimes_R \Bigl(\Gamma\bigl(C_R, \cap\bigr)\big/\mathrm{range}(\mathrm{diff}_R)\Bigr)
    \;\cong\;
    \Gamma\bigl(C_{R'}, \cap\bigr)\big/\mathrm{range}(\mathrm{diff}_{R'}),
    \qquad r' \otimes [s] \longmapsto [\,r' \cdot \rho^{\sharp} s\,],
  \]
  where \(\cap\) denotes the overlap of the respective base-changed covers
  (\ref{def:TwoCover_H1Cok}). No hypothesis on \(R \to R'\) is needed.
\end{theorem}
\begin{proof}
  \leanok
  Base change commutes with cokernels (\ref{thm:quot_range_baseChange}, with
  \(N = \Gamma(C_R, \cap)\), \(f = \mathrm{diff}_R\), \(S = R'\)):
  \(R' \otimes_R \bigl(N/\mathrm{range}\,f\bigr) \cong (R' \otimes_R N)/
  \mathrm{range}(f \otimes_R R')\). Then pass along the overlap term equivalence, which
  identifies \(\mathrm{range}(\mathrm{diff}_R \otimes_R R')\) with
  \(\mathrm{range}(\mathrm{diff}_{R'})\) (\ref{lem:relDiffBaseChange}), hence descends to
  the quotients. The composite sends \(r' \otimes [s]\) to
  \([\,r' \cdot \rho^{\sharp} s\,]\) by \ref{thm:relTermBaseChange}.
\end{proof}

\begin{theorem}[CBC-2: \(H^1\) of the relative curve commutes with change of the test ring]
  \label{thm:relH1BaseChange}
  \lean{AlgebraicGeometry.relH1BaseChange}
  \uses{def:relTwoCoverH1, thm:relH1CokBaseChange}
  \leanok
  \source{kleiman-picard}
  For every homomorphism \(R \to R'\) of commutative \(k\)-algebras there is an
  \(R'\)-linear equivalence
  \[
    R' \otimes_R H^1\bigl(C_R, \struct{C_R/R}\bigr) \;\cong\;
      H^1\bigl(C_{R'}, \struct{C_{R'}/R'}\bigr) :
  \]
  degree-one cohomology of the structure sheaf of the relative curve commutes with
  \emph{every} base change of the test ring. No flatness of \(R \to R'\) is required:
  the two-cover \(H^1\) is the cokernel of a two-term complex whose terms base-change
  freely, and cokernels commute with arbitrary base change.
\end{theorem}
\begin{proof}
  \leanok
  Conjugate the cokernel equivalence of \ref{thm:relH1CokBaseChange} by the relative
  two-cover identifications \(H^1(C_R, \struct{C_R/R}) \cong
  \Gamma(C_R, \cap)/\mathrm{range}(\mathrm{diff}_R)\) of \ref{def:relTwoCoverH1}, over
  \(R\) and over \(R'\).
\end{proof}

\begin{definition}[Base change of the unit cocycle]
  \label{def:relCocycleBaseChange}
  \lean{AlgebraicGeometry.relCocycleBaseChange}
  \uses{def:relTwistSheaf, def:relSectionsMap, lem:relCurveMap_proj}
  \leanok
  For a unit cocycle \(g\) on the base-changed cover of \(C_R\), the image
  \(\rho^{\sharp}(g)\) is a unit cocycle on the base-changed cover of \(C_{R'}\): the
  ring homomorphism \(\rho^{\sharp}\) preserves units, and the overlaps correspond under
  \(\rho\) (\ref{lem:relCurveMap_proj}).
\end{definition}

\begin{definition}[CBC-1 for the twisted sheaf: term equivalences]
  \label{def:relTwistTermBaseChange}
  \lean{AlgebraicGeometry.relTwistTermBaseChange₀, AlgebraicGeometry.relTwistTermBaseChange₁,
    AlgebraicGeometry.relTwistOverlapBaseChange}
  \uses{def:relTwistSheaf, def:relCocycleBaseChange, def:twistTriv, thm:relTermBaseChange,
    lem:relDiffBaseChange}
  \leanok
  Let \(g\) be a unit cocycle on the base-changed cover of \(C_R\) and
  \(g' = \rho^{\sharp}(g)\) its base change (\ref{def:relCocycleBaseChange}). For each
  term of the twisted two-cover complex --- the two charts and their overlap --- the
  composite of the chart trivialisation over \(R\) (\ref{def:twistTriv}), the
  structure-sheaf term base change of \ref{thm:relTermBaseChange} (at the overlap: the
  codomain equivalence of \ref{lem:relDiffBaseChange}), and the inverse chart
  trivialisation over \(R'\) is an \(R'\)-linear equivalence
  \[
    R' \otimes_R F_g(\text{term over } R) \;\cong\; F_{g'}(\text{term over } R') .
  \]
  On the overlap, the first-chart trivialisation (at
  \(\mathrm{pr}_1^{-1}V_0 \cap \mathrm{pr}_1^{-1}V_1 \le \mathrm{pr}_1^{-1}V_0\)) is
  used.
\end{definition}
